\chapter{Kode Program}
\label{lamp:C}

\begin{lstlisting}[language=JavaScript, caption={index.html}]
<!DOCTYPE html>
<html lang="id">

<head>
  <meta charset="UTF-8">
  <meta name="viewport" content="width=device-width, initial-scale=1.0">
  <title>Rusunami The Jarrdin Cihampelas</title>
  <script src="https://cdn.jsdelivr.net/npm/chart.js"></script>
  <link rel="preconnect" href="https://fonts.googleapis.com">
  <link rel="preconnect" href="https://fonts.gstatic.com" crossorigin>
  <link href="https://fonts.googleapis.com/css2?family=Montserrat:wght@400;500;600;700&display=swap" rel="stylesheet">
  <style>
    :root {
      --theme: #399051;
      --themeSubtle: #f1f6f2;
      --themeHover: #2e8546;
      --themeTua: #0a7e66;
      --theme2: #daaf15;
      --theme2Hover: #caa007;
      --gradient-theme: linear-gradient(to right, var(--theme), var(--themeTua));
      --box-shadow: 0px 4px 25px 0px rgba(0, 0, 0, 0.06);
      --border: #d4dcff;
      --sidebar-width: 240px;
      --sidebar-collapsed-width: 64px;
    }

    * {
      font-family: "Montserrat", sans-serif !important;
      box-sizing: border-box;
    }

    body {
      margin: 0;
      padding: 0;
      background-color: #f3f7fb;
      display: flex;
      width: 100vw;
      height: 100vh;
      overflow: hidden;
    }

    /* -- APP SHELL -- */
    #app-shell {
      display: flex;
      width: 100%;
      height: 100%;
      position: fixed;
    }

    /* -- SIDEBAR -- */
    .sidebar {
      position: relative;
      width: var(--sidebar-width);
      min-width: var(--sidebar-width);
      height: 100%;
      background-color: var(--themeSubtle);
      display: flex;
      flex-direction: column;
      justify-content: space-between;
      padding: 16px 0;
      transition: width 0.3s ease, min-width 0.3s ease;
      overflow: hidden;
      z-index: 100;
    }

    .sidebar.collapsed {
      width: var(--sidebar-collapsed-width);
      min-width: var(--sidebar-collapsed-width);
    }

    .sidebar-head {
      padding: 0 16px;
      display: flex;
      flex-direction: column;
      gap: 16px;
    }

    .sidebar-toggle-row {
      display: flex;
      justify-content: space-between;
      align-items: center;
    }

    .sidebar-user-info {
      display: flex;
      flex-direction: column;
      overflow: hidden;
    }

    .sidebar-user-info p {
      margin: 0;
      white-space: nowrap;
      overflow: hidden;
      text-overflow: ellipsis;
      font-size: 13px;
    }

    .sidebar-user-info .user-name {
      font-weight: 600;
      font-size: 14px;
    }

    .sidebar.collapsed .sidebar-user-info {
      display: none;
    }

    /* sidebar nav menu */
    .sidebar-menu {
      list-style: none;
      margin: 16px 0 0 0;
      padding: 0;
      flex: 1;
      display: flex;
      flex-direction: column;
      gap: 4px;
      overflow-y: auto;
    }

    .sidebar-menu li {
      display: flex;
    }

    .menu-item {
      display: flex;
      align-items: center;
      gap: 12px;
      width: 100%;
      padding: 10px 16px;
      border: none;
      background-color: transparent;
      color: #333;
      cursor: pointer;
      font-size: 13px;
      font-weight: 500;
      text-align: left;
      text-decoration: none;
      border-radius: 0;
      transition: background-color 0.2s ease, color 0.2s ease;
      white-space: nowrap;
      overflow: hidden;
    }

    .menu-item .menu-icon {
      font-size: 18px;
      flex-shrink: 0;
      width: 22px;
      text-align: center;
    }

    .menu-item .menu-label {
      overflow: hidden;
      text-overflow: ellipsis;
    }

    .sidebar.collapsed .menu-item .menu-label {
      display: none;
    }

    .menu-item:hover {
      background-color: #ffffff;
      color: var(--theme);
    }

    .menu-item.active {
      background-color: #ffffff;
      color: var(--theme);
      font-weight: 600;
      border-right: 3px solid var(--theme);
    }

    .sidebar-footer {
      padding: 0 16px;
    }

    /* toggle button inside sidebar */
    .sidebar-toggle-btn {
      background-color: var(--theme);
      color: white;
      border: none;
      border-radius: 6px;
      padding: 6px 10px;
      cursor: pointer;
      font-size: 16px;
      flex-shrink: 0;
      transition: background-color 0.2s;
    }

    .sidebar-toggle-btn:hover {
      background-color: var(--themeHover);
    }

    .btn-logout {
      display: flex;
      align-items: center;
      gap: 8px;
      width: 100%;
      padding: 10px 16px;
      background-color: transparent;
      color: #dc3545;
      border: 1px solid #dc3545;
      border-radius: 6px;
      cursor: pointer;
      font-size: 13px;
      font-weight: 500;
      transition: background-color 0.2s, color 0.2s;
      white-space: nowrap;
      overflow: hidden;
    }

    .btn-logout:hover {
      background-color: #dc3545;
      color: white;
    }

    .sidebar.collapsed .btn-logout .menu-label {
      display: none;
    }

    /* -- MAIN AREA -- */
    .main-area {
      flex: 1;
      display: flex;
      flex-direction: column;
      overflow: hidden;
      height: 100%;
    }

    /* Header bar (green top bar) */
    .content-header {
      display: flex;
      align-items: center;
      justify-content: space-between;
      padding: 12px 24px;
      background-color: var(--theme);
      box-shadow: 0 2px 8px rgba(0,0,0,0.1);
      flex-shrink: 0;
    }

    .content-header h5 {
      margin: 0;
      color: white;
      font-weight: 500;
      font-size: 15px;
    }

    .header-actions {
      display: flex;
      align-items: center;
      gap: 12px;
    }

    .search-box {
      display: flex;
      align-items: center;
      background: rgba(255,255,255,0.15);
      border: 1px solid rgba(255,255,255,0.4);
      border-radius: 6px;
      padding: 5px 10px;
      gap: 6px;
    }

    .search-box input {
      background: transparent;
      border: none;
      outline: none;
      color: white;
      font-size: 13px;
      width: 180px;
      padding: 0;
    }

    .search-box input::placeholder {
      color: rgba(255,255,255,0.7);
    }

    .search-icon {
      color: rgba(255,255,255,0.8);
      font-size: 14px;
    }

    .btn-insert-header {
      background-color: var(--theme2);
      color: white;
      border: none;
      border-radius: 6px;
      padding: 7px 14px;
      cursor: pointer;
      font-size: 13px;
      font-weight: 500;
      transition: background-color 0.2s;
      white-space: nowrap;
    }

    .btn-insert-header:hover {
      background-color: var(--theme2Hover);
    }

    /* divider line below header */
    .header-divider {
      height: 1px;
      background-color: #ACACAC;
      flex-shrink: 0;
    }

    /* Scrollable content area */
    .content-body {
      flex: 1;
      overflow-y: auto;
      padding: 24px;
      background-color: white;
    }

    /* -- PAGE TITLE AREA -- */
    .page-title-bar {
      display: flex;
      align-items: center;
      justify-content: space-between;
      margin-bottom: 20px;
      padding-bottom: 16px;
      border-bottom: 1px solid #e8e8e8;
    }

    .page-title-bar h2 {
      margin: 0;
      font-size: 18px;
      font-weight: 600;
      color: #333;
    }

    /* -- FORMS -- */
    .container {
      max-width: 100%;
      margin: 0;
      background-color: white;
      padding: 0;
    }

    /* The old top header (logo + title) shown only on login page - hidden in app */
    .app-top-header {
      display: flex;
      align-items: center;
      gap: 16px;
      margin-bottom: 24px;
      border-bottom: 3px solid var(--theme);
      padding-bottom: 16px;
    }

    .logo {
      width: 60px;
      height: 60px;
      border-radius: 50%;
      object-fit: cover;
    }

    .title {
      margin: 0;
      color: var(--theme);
      font-size: 20px;
      font-weight: 700;
    }

    .form-group {
      margin-bottom: 15px;
    }

    label {
      display: block;
      margin-bottom: 5px;
      font-weight: 600;
      font-size: 13px;
      color: #444;
    }

    input,
    select,
    textarea {
      width: 100%;
      padding: 8px 10px;
      border: 1px solid #d9d9d9;
      border-radius: 6px;
      font-size: 13px;
      transition: border-color 0.2s;
    }

    input:focus,
    select:focus,
    textarea:focus {
      border-color: var(--theme);
      outline: none;
      box-shadow: 0 0 0 2px rgba(57, 144, 81, 0.15);
    }

    button {
      background-color: var(--theme);
      color: white;
      padding: 8px 16px;
      border: none;
      border-radius: 6px;
      cursor: pointer;
      margin: 4px;
      font-size: 13px;
      font-weight: 500;
      transition: background-color 0.2s;
    }

    button:hover {
      background-color: var(--themeHover);
    }

    .btn-danger {
      background-color: #dc3545 !important;
      color: white !important;
    }

    .btn-danger:hover {
      background-color: #c82333 !important;
    }

    .btn-success {
      background-color: var(--theme) !important;
      color: white !important;
    }

    .btn-success:hover {
      background-color: var(--themeHover) !important;
    }

    .alert {
      padding: 12px 16px;
      margin-bottom: 16px;
      border: 1px solid transparent;
      border-radius: 6px;
      font-size: 13px;
    }

    .alert-success {
      color: #155724;
      background-color: #d4edda;
      border-color: #c3e6cb;
    }

    .alert-danger {
      color: #721c24;
      background-color: #f8d7da;
      border-color: #f5c6cb;
    }

    .alert-info {
      color: #0c5460;
      background-color: #d1ecf1;
      border-color: #bee5eb;
    }

    /* -- TABLE -- */
    .table {
      width: 100%;
      border-collapse: collapse;
      margin-top: 12px;
      font-size: 13px;
    }

    .table th,
    .table td {
      border: 1px solid #f0f0f0;
      padding: 10px 12px;
      text-align: left;
      white-space: nowrap;
      vertical-align: middle;
    }

    .table th {
      background-color: var(--themeSubtle);
      color: #333;
      font-weight: 600;
    }

    .table tbody tr:hover {
      background-color: #fafafa;
    }

    .hidden {
      display: none;
    }

    /* removed old menu-title / menu-box from sidebar - not needed */
    .menu-title { display: none; }
    .menu-box { display: none; }

    .admin-only {
      display: none !important;
    }

    .dashboard-grid {
      display: grid;
      grid-template-columns: repeat(auto-fit, minmax(420px, 1fr));
      gap: 20px;
    }

    .dashboard-card {
      background-color: #fafafa;
      border: 1px solid #e8e8e8;
      border-radius: 8px;
      padding: 18px;
    }

    /* -- PAGINATION -- */
    .pagination-container {
      margin: 16px 0;
      display: flex;
      align-items: center;
      justify-content: center;
      gap: 4px;
      flex-wrap: wrap;
    }

    .pagination-container button {
      margin: 2px;
      padding: 6px 12px;
      font-size: 13px;
      background-color: white;
      color: #333;
      border: 1px solid #d9d9d9;
      border-radius: 6px;
    }

    .pagination-container button:hover {
      border-color: var(--theme);
      color: var(--theme);
      background-color: white;
    }

    .pagination-container button.active {
      background-color: var(--theme);
      color: white;
      border-color: var(--theme);
      font-weight: 600;
    }

    .pagination-container button:disabled {
      background-color: #f5f5f5;
      color: #bbb;
      border-color: #d9d9d9;
      cursor: not-allowed;
      opacity: 0.8;
    }

    .pagination-container .pagination-info {
      margin-left: 12px;
      font-size: 13px;
      color: #666;
    }

    .pagination-container span.page-ellipsis {
      display: inline-flex;
      align-items: center;
      padding: 6px 8px;
      margin: 2px;
      color: #666;
    }

    /* -- SCROLLBAR -- */
    .content-body::-webkit-scrollbar { width: 6px; }
    .content-body::-webkit-scrollbar-track { background: transparent; }
    .content-body::-webkit-scrollbar-thumb { background-color: #ccc; border-radius: 4px; }
    .content-body::-webkit-scrollbar-thumb:hover { background-color: #aaa; }

    /* -- PAGE CONTENT (inner) -- */
    .page-content { display: none; }
    .page-content.active-page { display: block; }
  </style>
</head>

<body>

  <!-- -- APP SHELL (shown after login) -- -->
  <div id="app-shell" style="display:none;">
    <!-- SIDEBAR -->
    <div class="sidebar" id="sidebar">
      <div>
        <div class="sidebar-head">
          <div class="sidebar-toggle-row">
            <div class="sidebar-user-info" id="sidebar-user-info">
              <p>Selamat datang,</p>
              <p class="user-name" id="sidebar-user-name">...</p>
              <p class="user-role" id="sidebar-user-role" style="display:none;"></p>
            </div>
            <button class="sidebar-toggle-btn" id="sidebar-toggle-btn" onclick="toggleSidebar()">&#9776;</button>
          </div>
        </div>
        <ul class="sidebar-menu">
          <!-- <li><a class="menu-item" href="https://member.thejarrdin.com" target="_blank"><span class="menu-icon">??</span><span class="menu-label">Home</span></a></li> -->
          <li><button class="menu-item" id="menu-dashboard" onclick="showPage('dashboard-section')"><span class="menu-icon">??</span><span class="menu-label">Dashboard</span></button></li>
          <li><button class="menu-item" id="menu-input" onclick="showPage('input-form-section')"><span class="menu-icon">??</span><span class="menu-label">Input Data Penyewa</span></button></li>
          <li><button class="menu-item" id="menu-view" onclick="showPage('view-data-section')"><span class="menu-icon">??</span><span class="menu-label">Lihat / Edit Data</span></button></li>
          <li><button class="menu-item admin-only" id="menu-units" onclick="showPage('unit-management-section')"><span class="menu-icon">??</span><span class="menu-label">Manajemen Unit</span></button></li>
        </ul>
      </div>
      <div class="sidebar-footer">
        <button class="btn-logout" onclick="logout()">
          <span class="menu-icon" style="font-size: 16px;">??</span>
          <span class="menu-label">Keluar</span>
        </button>
      </div>
    </div>

    <!-- MAIN AREA -->
    <div class="main-area">
      <!-- GREEN HEADER -->
      <div class="content-header" id="content-header">
        <h5 id="page-title-display">Penyewaan Unit</h5>
        <div class="header-actions">
          <div class="search-box" id="header-search-box" style="display:none;">
            <span class="search-icon">??</span>
            <input type="text" id="header-search-input" placeholder="Cari..." oninput="onHeaderSearch(this.value)" />
          </div>
          <button class="btn-insert-header" id="header-insert-btn" style="display:none;" onclick="showPage('input-form-section')">
            + Tambah Data Penyewa
          </button>
        </div>
      </div>
      <!-- DIVIDER -->
      <div class="header-divider"></div>

      <!-- SCROLLABLE CONTENT -->
      <div class="content-body">

        <!-- INPUT FORM -->
        <div id="input-form-section" class="page-content">
          <div class="page-title-bar">
            <h2>?? Form Input Data Penyewa</h2>
          </div>
          <form id="rental-input-form">
            <div class="form-group"><label>Nama Penyewa</label><input type="text" id="nama" required></div>
            <div class="form-group"><label>NIK</label><input type="text" id="nik_input" pattern="\d{16}" minlength="16"
              maxlength="16" inputmode="numeric" oninput="this.value = this.value.replace(/\D/g, '').slice(0,16)"
              title="NIK harus 16 digit angka" required></div>
            <div class="form-group"><label>Status Pasutri</label>
              <select id="status_pasutri_input">
                <option value="Bukan">Bukan</option>
                <option value="Ya">Ya</option>
              </select>
            </div>
            <div class="form-group"><label>Pilih Unit</label><select id="unit-select" onchange="populateUnitFields()" required>
                <option value="">Memuat unit...</option>
              </select></div>
            <div class="form-group hidden"><label>Tower</label><input type="text" id="tower" readonly></div>
            <div class="form-group hidden"><label>Lantai</label><input type="text" id="lantai" readonly></div>
            <div class="form-group hidden"><label>Nomor Unit</label><input type="text" id="unit" readonly></div>
            <div class="form-group"><label>Kewarganegaraan</label><select id="status_kewarganegaraan">
                <option>WNI</option>
                <option>WNA</option>
              </select></div>
            <div class="form-group"><label>Jenis Sewa</label><select id="jenis_sewa" onchange="handleJenisSewaChange()">
                <option>Harian</option>
                <option>Mingguan</option>
                <option>Bulanan</option>
              </select></div>
            <div class="form-group hidden" id="bulanan-options"><label for="jumlah_bulan">Jumlah Bulan</label><input
                type="number" id="jumlah_bulan" min="1" value="1" oninput="calculateCheckoutAndDuration()"></div>
            <div class="form-group"><label>Metode Pembayaran</label><select id="metode_pembayaran"
                onchange="toggleMetodeLain()">
                <option>Cash</option>
                <option>Kartu Kredit</option>
                <option>Kartu Debit</option>
                <option>QRIS</option>
                <option value="Others">Others</option>
              </select></div>
            <div class="form-group hidden" id="metode-lain-group"><label>Metode Lainnya</label><input type="text"
                id="metode_lain"></div>
            <div class="form-group"><label>Tanggal Check-In</label><input type="date" id="tanggal_checkin"
                onchange="calculateCheckoutAndDuration()" required></div>
            <div class="form-group"><label>Waktu Check-In</label><input type="time" id="waktu_checkin" required></div>
            <div class="form-group"><label>Tanggal Check-Out</label><input type="date" id="tanggal_checkout"
                onchange="calculateCheckoutAndDuration()" required></div>
            <div class="alert alert-info" id="lama-menginap-info">Lama Menginap: 1 Hari</div>
            <button type="button" class="btn-success" onclick="saveRental()">Simpan Data</button>
          </form>
        </div>

        <!-- VIEW / EDIT DATA -->
        <div id="view-data-section" class="page-content">
          <div class="page-title-bar">
            <h2>?? Data Penyewa</h2>
            <div style="display:flex; gap:8px; flex-wrap:wrap;">
              <button class="btn-success" onclick="downloadRentalsCSV()" style="margin:0;">?? Download Data Penyewa</button>
              <button class="btn-success" onclick="downloadAllLogsCSV()" style="margin:0;">?? Download Semua Log</button>
              <button class="admin-only" onclick="checkDuplicates()" style="margin:0;">?? Cek Duplikat</button>
              <button onclick="loadRentalsData('default',true)" style="margin:0;">?? Refresh</button>
            </div>
          </div>
          <div class="admin-only" style="display:flex; align-items:center; gap:8px; margin-bottom:12px;">
            <input type="checkbox" id="comment-filter" onchange="loadRentalsData('default', true)" style="width:auto; margin:0;">
            <label for="comment-filter" style="margin:0; font-weight:normal; white-space:nowrap; font-size:13px;">Tampilkan data dengan pelanggaran saja</label>
          </div>
          <div id="rentals-table" style="overflow-x: auto;"></div>
          <hr style="margin:16px 0; border:none; border-top:1px solid #e8e8e8;">
          <h3 style="font-size:15px; margin-bottom:12px;">?? Edit Data Penyewa</h3>
          <div id="rentals-expanders"></div>
          <div id="pagination-controls-bottom" class="pagination-container"></div>
        </div>

        <!-- DASHBOARD -->
        <div id="dashboard-section" class="page-content">
          <div class="page-title-bar">
            <h2>?? Dashboard</h2>
          </div>
          <div class="dashboard-grid">
            <div class="dashboard-card">
              <h3 style="font-size:15px; font-weight:600; margin-bottom:12px; border-bottom:2px solid var(--themeSubtle); padding-bottom:8px;">Statistik Penyewaan Unit</h3>
              <div style="margin-bottom: 12px;">
                <label for="stats-filter" style="font-size:13px;">Filter Periode:</label>
                <select id="stats-filter" onchange="loadDashboardStats()" style="width:auto; display:inline-block; margin-left:8px;">
                  <option value="all">Semua Waktu</option>
                  <option value="monthly">Bulanan</option>
                  <option value="weekly">Mingguan</option>
                </select>
              </div>
              <h4 style="text-align:center; font-weight:500; font-size:13px; margin-bottom:8px;">Jumlah Unit Disewakan per Periode</h4>
              <canvas id="rentedUnitsChart"></canvas>
              <hr style="margin:16px 0; border:none; border-top:1px solid #e8e8e8;">
              <h4 style="text-align:center; font-weight:500; font-size:13px; margin-bottom:8px;">Top 5 Unit Paling Populer</h4>
              <canvas id="topUnitsChart"></canvas>
              <div id="admin-stats-container" class="admin-only" style="margin-top:16px; display:none;">
                <hr style="margin:16px 0; border:none; border-top:1px solid #e8e8e8;">
                <h3 style="font-size:15px; font-weight:600; margin-bottom:8px;">Statistik Detail Setiap Pelaku Komersil</h3>
                <h4 style="text-align:center; font-weight:500; font-size:13px; margin-bottom:8px;">Jumlah Unit Disewakan (per Pelaku Komersil)</h4>
                <canvas id="rentedUnitsByAgentChart"></canvas>
                <hr style="margin:16px 0; border:none; border-top:1px solid #e8e8e8;">
                <h4 style="text-align:center; font-weight:500; font-size:13px; margin-bottom:8px;">Top 5 Unit Populer (per Pelaku Komersil)</h4>
                <canvas id="topUnitsByAgentChart"></canvas>
              </div>
            </div>
            <div class="dashboard-card">
              <h3 style="font-size:15px; font-weight:600; margin-bottom:12px; border-bottom:2px solid var(--themeSubtle); padding-bottom:8px;">Laporan Pelanggaran</h3>
              <div id="violations-container" style="max-height:400px; overflow-y:auto;"></div>
              <hr style="margin:16px 0; border:none; border-top:1px solid #e8e8e8;">
              <div class="admin-only">
                <h3 style="font-size:15px; font-weight:600; margin-bottom:12px;">Upload Laporan Baru</h3>
                <div class="form-group"><label>Pilih Penyewa</label><select id="violation-rental-id"></select></div>
                <div class="form-group"><label>Upload Foto</label><input type="file" id="violation-photo-file" accept="image/*"></div>
                <div class="form-group"><label>Deskripsi</label><input type="text" id="violation-desc"></div>
                <button class="btn-success" onclick="submitViolation()">Upload Laporan</button>
              </div>
            </div>
          </div>
        </div>

        <!-- UNIT MANAGEMENT -->
        <div id="unit-management-section" class="page-content">
          <div class="page-title-bar">
            <h2>?? Manajemen Unit Pelaku Komersil</h2>
          </div>
          <div class="form-group">
            <label>Pilih Pengguna Pelaku Komersil</label>
            <select id="manage-agent-select" onchange="loadUnitsForAgent()"></select>
          </div>
          <div id="unit-list-container"></div>
          <hr style="margin:16px 0; border:none; border-top:1px solid #e8e8e8;">
          <h3 style="font-size:15px; font-weight:600; margin-bottom:12px;">Tambah Unit Baru</h3>
          <div class="form-group">
            <label>Tower</label>
            <select id="new-unit-tower">
              <option value="A">A</option>
              <option value="B">B</option>
              <option value="C">C</option>
              <option value="D">D</option>
            </select>
          </div>
          <div class="form-group">
            <label>Lantai</label>
            <input type="text" id="new-unit-lantai">
          </div>
          <div class="form-group">
            <label>Nomor</label>
            <input type="text" id="new-unit-nomor">
          </div>
          <button class="btn-success" onclick="addUnit()" style="width:100%; margin:10px 0 0 0;">Tambah Unit</button>
        </div>

      </div><!-- end content-body -->
    </div><!-- end main-area -->
  </div><!-- end app-shell -->

  <div id="alert-container" style="position:fixed; top:20px; right:20px; z-index:9999;"></div>

  <div id="log-modal"
    style="display:none; position:fixed; z-index:2000; left:0; top:0; width:100%; height:100%; overflow:auto; background-color:rgba(0,0,0,0.5);">
    <div style="background-color:#fefefe; margin:10% auto; padding:20px; border:1px solid #e8e8e8; width:80%; max-width:900px; border-radius:8px; box-shadow:0 4px 24px rgba(0,0,0,0.12);">
      <span style="color:#aaa; float:right; font-size:28px; font-weight:bold; cursor:pointer;"
        onclick="closeLogModal()">&times;</span>
      <div style="display:flex; justify-content:space-between; align-items:center;">
        <h2 id="log-modal-title" style="font-size:16px; font-weight:600;">Riwayat Perubahan</h2>
        <button class="btn-success" id="download-log-tenant-csv" onclick="" style="margin:0;">?? Download Log Ini (CSV)</button>
      </div>
      <div id="log-modal-body" style="max-height:400px; overflow-y:auto; margin-top:15px;"></div>
    </div>
  </div>

  <script>
    let currentUser = null;
    let rentalsData = [];

    let currentPage = 1;
    const itemsPerPage = 10;
    let currentViewMode = 'default';
    // Chart variables
    let rentedUnitsChart, topUnitsChart, rentedUnitsByAgentChart, topUnitsByAgentChart;

    function getAuthHeaders(extra = {}) {
      const session = JSON.parse(localStorage.getItem('userSession') || 'null');
      const token = session?.token || '';
      return { 'Authorization': 'Bearer ' + token, ...extra };
    }

    function isUserAdmin(user) {
      if (!user) return false;
      // Normalize: handle both server-normalized (roles, string[]) and raw SRusun format (Role, object[])
      let rolesArr = user.roles;
      if (!rolesArr && Array.isArray(user.Role)) {
        rolesArr = user.Role.map(r => (typeof r === 'object' ? r.Nama : r));
      }
      if (!rolesArr || rolesArr.length === 0) return false;
      // SRusun roles: 'Pengelola' = admin, 'Pelaku Komersil' = regular user
      return rolesArr.some(r => (typeof r === 'object' ? r.Nama : r) === 'Pengelola');
    }


    document.addEventListener('DOMContentLoaded', () => {
      checkAuthStatus();
    });

    function toggleSidebar() {
      document.getElementById('sidebar').classList.toggle('collapsed');
    }

    async function checkAuthStatus() {
      // Accept token from OTP redirect via URL param (?token=xxx)
      const urlParams = new URLSearchParams(window.location.search);

      // Handle logout param to remove session
      if (urlParams.get('logout') === 'true') {
        localStorage.removeItem('userSession');
        currentUser = null;
        window.history.replaceState({}, '', window.location.pathname);
        window.location.href = '/login.html'; // window.location.href = '/login.html'; // window.location.href = 'https://member.thejarrdin.com/login?from=visitor';
        return;
      }

      const urlToken = urlParams.get('token');
      if (urlToken) {
        try {
          const decodedToken = decodeURIComponent(atob(urlToken));
          // decodedToken is the raw JWT; decode its payload for the object
          const payloadB64 = decodedToken.split('.')[1];
          const payloadObj = payloadB64 ? JSON.parse(atob(payloadB64)) : {};
          localStorage.setItem('userSession', JSON.stringify({ token: decodedToken, object: payloadObj }));
        } catch (e) {
          console.error('Failed to decode token param', e);
        }
        window.history.replaceState({}, '', window.location.pathname);
      }

      // Accept session forwarded from SRusun-WebMember (?session=base64encodedJSON)
      const sessionParam = urlParams.get('session');
      if (sessionParam) {
        try {
          const decoded = decodeURIComponent(atob(sessionParam));
          const sessionObj = JSON.parse(decoded);
          // Support both SRusun format { AuthKey, dataUser } and plain JWT string
          const authToken = sessionObj?.AuthKey || (typeof sessionObj === 'string' ? sessionObj : null);
          if (authToken) {
            localStorage.setItem('userSession', JSON.stringify({ token: authToken, object: sessionObj }));
          }
          window.history.replaceState({}, '', window.location.pathname);
        } catch (e) {
          console.error('Failed to parse session param', e);
        }
      }

      // Also handle if localStorage was set directly with SRusun { AuthKey, dataUser } format
      let session = JSON.parse(localStorage.getItem('userSession') || 'null');
      if (session && session.AuthKey && !session.token) {
        session = { token: session.AuthKey, object: session };
        localStorage.setItem('userSession', JSON.stringify(session));
      }

      const token = session?.token;
      if (!token) {
        // Redirect to SRusun for auth, which will forward the session back via ?session= param
        window.location.href = '/login.html'; // window.location.href = '/login.html'; // window.location.href = 'https://member.thejarrdin.com/login?from=visitor';
        return;
      }
      try {
        const response = await fetch('/api/user', { headers: getAuthHeaders() });
        const data = await response.json();
        if (data.user) {
          currentUser = data.user;
          showMainSection();
        } else {
          localStorage.removeItem('userSession');
          window.location.href = '/login.html'; // window.location.href = '/login.html'; // window.location.href = 'https://member.thejarrdin.com/login?from=visitor';
        }
      } catch (error) {
        showAlert('Gagal terhubung ke server.', 'danger');
      }
    }

    function initializeDateTime() {
      const now = new Date();
      const today = now.toISOString().split('T')[0];
      const currentTime = now.toTimeString().slice(0, 5);
      document.getElementById('tanggal_checkin').value = today;
      document.getElementById('waktu_checkin').value = currentTime;
      document.getElementById('tanggal_checkout').value = today;
      calculateCheckoutAndDuration();
    }

    async function logout() {
      await fetch('/api/logout', { method: 'POST', headers: getAuthHeaders() });
      localStorage.removeItem('userSession');
      currentUser = null;
      window.location.href = '/login.html'; // window.location.href = '/login.html'; // window.location.href = 'https://member.thejarrdin.com/logout';
    }

    // Page titles and header options per page
    const PAGE_META = {
      'dashboard-section':       { title: 'Dashboard',           search: false, insert: false },
      'input-form-section':      { title: 'Input Data Penyewa',  search: false, insert: false },
      'view-data-section':       { title: 'Lihat / Edit Data',   search: true,  insert: false },
      'unit-management-section': { title: 'Manajemen Unit',      search: false, insert: false },
    };

    function showPage(pageId) {
      // Guard: Pelaku Komersil cannot access admin-only pages
      if (pageId === 'unit-management-section' && !isUserAdmin(currentUser)) {
        showAlert('Akses ditolak. Halaman ini hanya untuk Pengelola.', 'danger');
        return;
      }
      document.querySelectorAll('.page-content').forEach(p => p.classList.remove('active-page'));
      document.getElementById(pageId).classList.add('active-page');

      // Update green header title
      const meta = PAGE_META[pageId] || { title: 'Penyewaan Unit', search: false, insert: false };
      const titleEl = document.getElementById('page-title-display');
      if (titleEl) titleEl.textContent = meta.title;

      // Show/hide search and insert button in header
      const searchBox = document.getElementById('header-search-box');
      const insertBtn = document.getElementById('header-insert-btn');
      const searchInput = document.getElementById('header-search-input');
      if (searchBox) { searchBox.style.display = meta.search ? 'flex' : 'none'; }
      if (searchInput) searchInput.value = '';
      if (insertBtn) { insertBtn.style.display = (meta.insert && isUserAdmin(currentUser)) ? '' : 'none'; }

      document.querySelectorAll('.menu-item').forEach(m => m.classList.remove('active'));
      let menuBtnId = `menu-${pageId.split('-')[0]}`;
      if (pageId === 'unit-management-section') {
        menuBtnId = 'menu-units';
      }
      const activeMenu = document.getElementById(menuBtnId);

      if (activeMenu) activeMenu.classList.add('active');

      if (pageId === 'input-form-section') {
        document.getElementById('rental-input-form').reset();
        initializeDateTime();
        loadMyUnitsForInput();
      }
      if (pageId === 'view-data-section') loadRentalsData('default', true);
      if (pageId === 'dashboard-section') loadDashboardData();
      if (pageId === 'unit-management-section') loadAgentsForManagement();
    }

    let searchTimeout = null;
    // Called by the header search input
    function onHeaderSearch(val) {
      if (searchTimeout) clearTimeout(searchTimeout);
      searchTimeout = setTimeout(() => {
        // Reset ke halaman 1 dan muat ulang data dari server
        loadRentalsData(currentViewMode, true);
      }, 500);
    }

    // Filter visible table rows by a text query (client-side, on the rendered table)
    function filterRentalsTable(query) {
      const q = (query || '').toLowerCase().trim();
      const table = document.querySelector('#rentals-table table');
      if (!table) return;
      const rows = table.querySelectorAll('tbody tr');
      rows.forEach(row => {
        row.style.display = (!q || row.textContent.toLowerCase().includes(q)) ? '' : 'none';
      });
    }

    function showMainSection() {
      const shell = document.getElementById('app-shell');
      shell.style.display = 'flex';

      // populate user info in sidebar
      const displayName = currentUser.nama || currentUser.Nama || currentUser.email || 'Pengguna';
      const userNameEl = document.getElementById('sidebar-user-name');
      if (userNameEl) userNameEl.textContent = displayName;

      // Role display is intentionally hidden in the sidebar

      const isAdmin = isUserAdmin(currentUser);
      if (isAdmin) {
        // Remove admin-only class so natural CSS (flex, block, etc.) restores correctly.
        // !important on the class means we must remove the class, not fight it with inline style.
        document.querySelectorAll('.admin-only').forEach(el => el.classList.remove('admin-only'));
      }
      // Non-admin: CSS .admin-only { display:none !important } already hides everything. No action needed.

      showPage('dashboard-section');
    }



    function toggleMetodeLain() {
      document.getElementById('metode-lain-group').classList.toggle('hidden', document.getElementById('metode_pembayaran').value !== 'Others');
    }

    function handleJenisSewaChange() {
      const jenisSewa = document.getElementById('jenis_sewa').value;
      document.getElementById('bulanan-options').classList.toggle('hidden', jenisSewa !== 'Bulanan');
      calculateCheckoutAndDuration();
    }

    function calculateCheckoutAndDuration() {
      try {
        const checkinInput = document.getElementById('tanggal_checkin');
        const checkoutEl = document.getElementById('tanggal_checkout');
        if (!checkinInput.value) return;

        const checkin = new Date(checkinInput.value);

        if (document.getElementById('jenis_sewa').value === 'Bulanan') {
          const months = parseInt(document.getElementById('jumlah_bulan').value) || 1;
          const newCheckout = new Date(checkin);
          newCheckout.setMonth(checkin.getMonth() + months);
          checkoutEl.valueAsDate = newCheckout;
        }

        if (!checkoutEl.value) return;

        const checkout = new Date(checkoutEl.value);
        const diffTime = Math.max(0, checkout - checkin);
        const diffDays = Math.ceil(diffTime / (1000 * 60 * 60 * 24));
        document.getElementById('lama-menginap-info').textContent = `Lama Menginap: ${diffDays} Hari`;
      } catch (e) {
        // handle errors
      }
    }

    async function saveRental() {
      // 1. Gather all data from the form fields
      const nikInput = document.getElementById('nik_input');
      const nikValue = nikInput.value.replace(/\D/g, '').slice(0, 16);
      nikInput.value = nikValue;

      const rentalData = {
        nama: document.getElementById('nama').value,
        nik: nikValue,
        status_pasutri: document.getElementById('status_pasutri_input').value,
        unit_number: document.getElementById('unit-select').value,
        status_kewarganegaraan: document.getElementById('status_kewarganegaraan').value,
        jenis_sewa: document.getElementById('jenis_sewa').value,
        metode_pembayaran: document.getElementById('metode_pembayaran').value,
        metode_lain: document.getElementById('metode_lain').value,
        tanggal_checkin: document.getElementById('tanggal_checkin').value,
        waktu_checkin: document.getElementById('waktu_checkin').value,
        tanggal_checkout: document.getElementById('tanggal_checkout').value
      };

      // 2. Perform basic validation
      console.log(rentalData)
      if (!rentalData.nama || !rentalData.nik || !rentalData.unit_number) {
        showAlert('Nama, NIK, dan Unit harus diisi.', 'danger');
        return;
      }
      if (!/^\d{16}$/.test(rentalData.nik)) {
        showAlert('NIK harus terdiri dari 16 digit angka.', 'danger');
        return;
      }

      // 3. Send the data to the backend API
      try {
        const response = await fetch('/api/rentals', {
          method: 'POST',
          headers: getAuthHeaders({ 'Content-Type': 'application/json' }),
          body: JSON.stringify(rentalData)
        });

        const result = await response.json();

        // 4. Show only the prediction warning when the model flags a violation risk
        if (!result.success) {
          showAlert(result.message, 'danger');
          return;
        }

        if (result.prediction && result.prediction.class === 1) {
          showPredictionNotification(result.prediction.notification || '??Penyewa Berpotensi Melakukan Pelanggaran!');
        }

        if (result.success) {
          document.getElementById('rental-input-form').reset();
          initializeDateTime(); // Reset date fields
          showPage('view-data-section'); // Go to the data table page
        }
      } catch (error) {
        console.error('Error saving rental data:', error);
        showAlert('Terjadi kesalahan koneksi saat menyimpan data.', 'danger');
      }
    }

    async function loadRentalsData(mode = 'default', isRefresh = false) {
      if (isRefresh) {
        currentPage = 1;
      }

      // Simpan mode saat ini untuk digunakan oleh tombol 'changePage'
      currentViewMode = mode;

      const tableContainer = document.getElementById('rentals-table');
      const expandersContainer = document.getElementById('rentals-expanders');
      const paginationBottom = document.getElementById('pagination-controls-bottom');

      tableContainer.innerHTML = '<p>Memuat data...</p>';
      expandersContainer.innerHTML = '';
      if (paginationBottom) paginationBottom.innerHTML = '';

      // --- 1. Buat URL berdasarkan mode & filter ---
      let baseUrl;
      let queryParams = new URLSearchParams();
      queryParams.append('page', currentPage);
      queryParams.append('limit', itemsPerPage);

      if (mode === 'default') {
        baseUrl = '/api/rentals';
        const commentFilter = document.getElementById('comment-filter').checked;
        queryParams.append('commentFilter', commentFilter);
        const searchInput = document.getElementById('header-search-input');
        if (searchInput && searchInput.value) {
          queryParams.append('search', searchInput.value.trim());
        }
      } else if (mode === 'duplicates') {
        baseUrl = '/api/rentals/duplicates';
      } else {
        tableContainer.innerHTML = '<p>Mode tampilan tidak valid.</p>';
        return;
      }

      const url = `${baseUrl}?${queryParams.toString()}`;

      // --- 2. Ambil data dari Server ---
      let responseData;
      let paginationData;

      try {
        const response = await fetch(url, { headers: getAuthHeaders() });
        const result = await response.json();

        if (!result.success) {
          showAlert(result.message, 'danger');
          tableContainer.innerHTML = `<div class="alert alert-danger">${result.message}</div>`;
          return;
        }

        responseData = result.data;
        paginationData = result.pagination;

        if (mode === 'duplicates' && responseData.length > 0) {
          showAlert(`Ditemukan ${paginationData.totalItems} data dengan NIK ganda.`, 'info');
        }

      } catch (error) {
        console.error('Error loading rental data:', error);
        showAlert('Gagal terhubung ke server untuk memuat data.', 'danger');
        tableContainer.innerHTML = '<div class="alert alert-danger">Gagal memuat data.</div>';
        return;
      }

      // --- 3. Render ---
      if (responseData.length === 0) {
        tableContainer.innerHTML = '<div class="alert alert-info">Tidak ada data.</div>';
        return;
      }

      // Render Paginasi (panggil SEBELUM render tabel/expander)
      renderPagination(paginationData.totalItems, paginationData.totalPages, paginationData.page);

      // --- 4. Render Tabel ---
      let tableHTML = `<table class="table"><thead><tr>
                <th>NIK</th><th>Nama</th><th>Kewarganegaraan</th><th>Status Perkawinan</th><th>Unit</th><th>Jenis Sewa</th>
                <th>Metode Pembayaran</th><th>Tgl Check In</th><th>Waktu Check In</th><th>Tgl Check Out</th><th>Waktu Checkout</th><th>Lama (Hari)</th>
                <th>Komentar</th><th>Email Pelaku Komersil</th><th>Diedit Oleh</th><th>Log</th>
            </tr></thead><tbody>`;

      responseData.forEach(r => {
        let paymentMethod = r.metode_pembayaran;
        if (paymentMethod === 'Others' && r.metode_lain) {
          paymentMethod = `Others (${r.metode_lain})`;
        }
        const checkinDate = r.tanggal_checkin ? new Date(r.tanggal_checkin).toLocaleDateString('id-ID') : '';
        const checkoutDate = r.tanggal_checkout ? new Date(r.tanggal_checkout).toLocaleDateString('id-ID') : '';

        tableHTML += `<tr>
                    <td>${r.nik || ''}</td>
                    <td>${r.nama || ''}</td>
                    <td>${r.status_kewarganegaraan || ''}</td>
                    <td>${r.status_pasutri || ''}</td>
                    <td>${r.unit_number || ''}</td>
                    <td>${r.jenis_sewa || ''}</td>
                    <td>${paymentMethod || ''}</td>
                    <td>${checkinDate}</td>
                    <td>${(r.waktu_checkin || '').substring(0,5)}</td>
                    <td>${checkoutDate}</td>
                    <td>${(r.waktu_checkout || '').substring(0,5)}</td>
                    <td>${r.lama_menginap ?? ''}</td>
                    <td>${(r.komentar || []).join(', ')}</td>
                    <td>${r.user_email || ''}</td>
                    <td>${r.diedit_oleh || ''}</td>
                    <td>
                        ${r.logs && r.logs.length > 0
            ? `<button onclick="showLogs(${r.id})">Detail Log</button>`
            : '<span>-</span>'}
                    </td>
                </tr>`;
      });
      tableHTML += '</tbody></table>';
      tableContainer.innerHTML = tableHTML;

      // --- 5. Render Expanders ---
      const pelanggaranList = await fetch('/api/pelanggaran-list', { headers: getAuthHeaders() }).then(res => res.json()).then(d => d.data);
      let expandersHTML = '';

      responseData.forEach((r) => {
        // ... (Logika pembuatan expanderHTML tetap sama seperti sebelumnya) ...
        expandersHTML += `<div class="expander" style="border: 1px solid #ddd; margin-bottom: 10px;">
                    <div class="expander-header" style="padding: 10px; background: #f0f0f0; cursor: pointer;" onclick="this.nextElementSibling.classList.toggle('hidden')">
                        ${r.nama} (Unit ${r.tower}-${r.lantai}-${r.unit})
                    </div>
                    <div class="expander-content hidden" style="padding: 15px;">
                        <div class="form-group"><label>Jenis Sewa</label><select id="edit_jenis_sewa_${r.id}">${['Harian', 'Mingguan', 'Bulanan'].map(o => `<option value="${o}" ${r.jenis_sewa === o ? 'selected' : ''}>${o}</option>`).join('')}</select></div>
                        <div class="form-group"><label>Metode Pembayaran</label><select id="edit_metode_pembayaran_${r.id}">${['Cash', 'Kartu Kredit', 'Kartu Debit', 'QRIS', 'Others'].map(o => `<option value="${o}" ${r.metode_pembayaran === o ? 'selected' : ''}>${o}</option>`).join('')}</select></div>
                        <div class="form-group"><label>Waktu Check-Out</label><input type="time" id="edit_waktu_checkout_${r.id}" value="${r.waktu_checkout || ''}"></div>
                        
                        <div class="form-group">
                            <label>Komentar</label>
                            <div>${pelanggaranList.map(p => {
          const checkboxId = `cb-${r.id}-${p.replace(/[^a-zA-Z0-9]/g, '-')}`;
          return `
                                    <div style="display: flex; align-items: center; margin-bottom: 5px;">
                                        <input type="checkbox" id="${checkboxId}" class="komentar-cb-${r.id}" value="${p}" ${(r.komentar || []).includes(p) ? 'checked' : ''} style="width: auto; margin-right: 8px;">
                                        <label for="${checkboxId}" style="font-weight: normal; margin-bottom: 0; display: inline;">${p}</label>
                                    </div>
                                `;
        }).join('')}</div>
                        </div>
                        <div style="text-align: left;">
                            <button class="btn-success" onclick="updateRental(${r.id})">Simpan Perubahan</button>
                        </div>
                    </div>
                </div>`;
      });
      expandersContainer.innerHTML = expandersHTML;
      rentalsData = responseData;
    }

    async function updateRental(id) {
      const data = {
        jenis_sewa: document.getElementById(`edit_jenis_sewa_${id}`).value,
        metode_pembayaran: document.getElementById(`edit_metode_pembayaran_${id}`).value,
        waktu_checkout: document.getElementById(`edit_waktu_checkout_${id}`).value,
        komentar: Array.from(document.querySelectorAll(`.komentar-cb-${id}:checked`)).map(cb => cb.value)
      };
      console.log(data)
      const response = await fetch(`/api/rentals/${id}`, {
        method: 'PUT',
        headers: getAuthHeaders({ 'Content-Type': 'application/json' }),
        body: JSON.stringify(data)
      });
      const result = await response.json();
      showAlert(result.message, result.success ? 'success' : 'danger');
      if (result.success) loadRentalsData();
    }

    async function checkDuplicates() {
      showAlert('Mencari data duplikat...', 'info');
      const response = await fetch('/api/rentals/duplicates', { headers: getAuthHeaders() });
      const result = await response.json();
      await loadRentalsData('duplicates', true);
      if (result.success) {
        if (result.data.length > 0) {
          showAlert(`Ditemukan ${result.data.length} data dengan NIK ganda. Menampilkan hasil.`, 'info');
          loadRentalsData(false, result.data);
        } else {
          showAlert('Tidak ditemukan data duplikat.', 'success');
        }
      }
    }

    function showLogs(rentalId) {
      const rental = rentalsData.find(r => r.id === rentalId);
      if (!rental || !rental.logs || rental.logs.length === 0) return;

      document.getElementById('log-modal-title').innerText = `Riwayat Perubahan: ${rental.nama}`;
      const modalBody = document.getElementById('log-modal-body');

      document.getElementById('download-log-tenant-csv').setAttribute('onclick', `downloadTenantLogCSV(${rental.id})`);

      let logsHTML = '<table class="table"><thead><tr><th>Waktu</th><th>Email</th><th>Field</th><th>Nilai Lama</th><th>Nilai Baru</th></tr></thead><tbody>';

      rental.logs.forEach(log => {
        logsHTML += `<tr>
                    <td>${new Date(log.timestamp).toLocaleString('id-ID')}</td>
                    <td>${log.email || ''}</td>
                    <td>${log.field_changed || ''}</td>
                    <td><div style="max-height: 50px; overflow-y: auto;">${log.old_value || ''}</div></td>
                    <td><div style="max-height: 50px; overflow-y: auto;">${log.new_value || ''}</div></td>
                </tr>`;
      });

      logsHTML += '</tbody></table>';
      modalBody.innerHTML = logsHTML;
      document.getElementById('log-modal').style.display = 'block';
    }

    function closeLogModal() {
      document.getElementById('log-modal').style.display = 'none';
    }

    window.onclick = function (event) {
      const modal = document.getElementById('log-modal');
      if (event.target == modal) {
        modal.style.display = "none";
      }
    }

    function downloadRentalsCSV() {
      const token = JSON.parse(localStorage.getItem('userSession') || 'null')?.token || '';
      window.location.href = `/api/rentals/csv?token=${token}`;
    }

    function downloadAllLogsCSV() {
      const token = JSON.parse(localStorage.getItem('userSession') || 'null')?.token || '';
      window.location.href = `/api/logs/csv?token=${token}`;
    }

    function downloadTenantLogCSV(rentalId) {
      if (!rentalId) return;
      const token = JSON.parse(localStorage.getItem('userSession') || 'null')?.token || '';
      window.location.href = `/api/rentals/${rentalId}/log/csv?token=${token}`;
    }

    function loadDashboardData() {
      loadDashboardStats();
      loadDashboardViolations();
      const isAdmin = isUserAdmin(currentUser);
      if (isAdmin) {
        loadRentalsForViolationDropdown();
      }
    }


    async function loadDashboardStats() {
      const filter = document.getElementById('stats-filter').value;
      const isAdmin = isUserAdmin(currentUser);
      const adminStatsContainer = document.getElementById('admin-stats-container');

      // Hancurkan chart lama agar bisa digambar ulang
      if (rentedUnitsChart) rentedUnitsChart.destroy();
      if (topUnitsChart) topUnitsChart.destroy();
      if (rentedUnitsByAgentChart) rentedUnitsByAgentChart.destroy();
      if (topUnitsByAgentChart) topUnitsByAgentChart.destroy();

      try {
        // --- 1. Ambil Data Rangkuman (Server-Side) ---
        // Kita panggil semua API rangkuman secara paralel agar lebih cepat
        const [rentalsResponse, popularResponse, agentPerfResponse] = await Promise.all([
          fetch(`/api/dashboard/rentals-over-time?filter=${filter}`, { headers: getAuthHeaders() }),
          fetch(`/api/dashboard/popular-units?filter=${filter}`, { headers: getAuthHeaders() }),
          isAdmin ? fetch(`/api/dashboard/agent-performance?filter=${filter}`, { headers: getAuthHeaders() }) : Promise.resolve(null)
        ]);

        // --- 2. Proses Chart 1: Jumlah Unit Disewakan (Line) ---
        const rentData = await rentalsResponse.json();
        if (rentData.success) {
          const rentedUnitsCanvas = document.getElementById('rentedUnitsChart');
          rentedUnitsChart = new Chart(rentedUnitsCanvas, {
            type: 'line',
            data: {
              labels: rentData.labels, // Langsung pakai data dari server
              datasets: [{
                label: 'Jumlah Unit Disewakan',
                data: rentData.data, // Langsung pakai data dari server
                borderColor: 'rgb(75, 192, 192)',
                tension: 0.1
              }]
            }
          });
        } else {
          showAlert('Gagal memuat statistik unit.', 'danger');
        }

        // --- 3. Proses Chart 2: Top 5 Unit Populer (Bar) ---
        const popData = await popularResponse.json();
        if (popData.success) {
          const topUnitsCanvas = document.getElementById('topUnitsChart');
          topUnitsChart = new Chart(topUnitsCanvas, {
            type: 'bar',
            data: {
              labels: popData.labels, // Langsung pakai data dari server
              datasets: [{
                label: 'Jumlah Sewa',
                data: popData.data, // Langsung pakai data dari server
                backgroundColor: 'rgba(255, 159, 64, 0.5)'
              }]
            },
            options: { indexAxis: 'y' }
          });
        } else {
          showAlert('Gagal memuat unit populer.', 'danger');
        }

        // --- 4. Proses Chart Admin ---
        if (isAdmin && agentPerfResponse) {
          adminStatsContainer.style.display = 'block';
          const agentPerfData = await agentPerfResponse.json();

          if (agentPerfData.success) {
            // Chart 3: Jumlah Sewa per Pelaku Komersil (Bar)
            // Kita gunakan data dari /agent-performance
            const agentLabels = agentPerfData.data.map(a => a.user_email);
            const agentTotals = agentPerfData.data.map(a => a.total_rented);

            const rentedUnitsByAgentCanvas = document.getElementById('rentedUnitsByAgentChart');
            rentedUnitsByAgentChart = new Chart(rentedUnitsByAgentCanvas, {
              type: 'bar', // Diubah jadi 'bar' agar lebih mudah dibaca
              data: {
                labels: agentLabels,
                datasets: [{
                  label: 'Total Unit Disewakan',
                  data: agentTotals,
                  backgroundColor: 'rgba(54, 162, 235, 0.5)'
                }]
              },
              options: { indexAxis: 'y' }
            });
          } else {
            showAlert('Gagal memuat performa Pelaku Komersil.', 'danger');
          }

          // Chart 4: Top 5 Unit (dengan info Pelaku Komersil)
          // Kita gunakan data dari /popular-units
          if (popData.success) {
            const topUnitsByAgentCanvas = document.getElementById('topUnitsByAgentChart');
            topUnitsByAgentChart = new Chart(topUnitsByAgentCanvas, {
              type: 'bar',
              data: {
                labels: popData.labels.map((label, index) => `${label} (${popData.agents[index] || 'N/A'})`), // Gabung label unit + Pelaku Komersil
                datasets: [{
                  label: 'Jumlah Sewa',
                  data: popData.data,
                  backgroundColor: 'rgba(153, 102, 255, 0.5)'
                }]
              },
              options: { indexAxis: 'y' }
            });
          }
        } else {
          adminStatsContainer.style.display = 'none';
        }
      } catch (error) {
        console.error("Error loading dashboard stats:", error);
        showAlert("Terjadi kesalahan saat memuat dashboard.", "danger");
      }
    }

    // Function to display violation photos directly and hide uploader for agents
    async function loadDashboardViolations() {
      const container = document.getElementById('violations-container');
      container.innerHTML = 'Memuat laporan...';
      const isAdmin = isUserAdmin(currentUser);

      const response = await fetch('/api/dashboard/violations', { headers: getAuthHeaders() });
      const result = await response.json();
      if (result.success) {
        let html = '';
        result.data.forEach(v => {
          const adminControls = isAdmin ? `
                        <div style="display: flex; gap: 10px; margin-top: 10px;">
                            <button class="btn-success" style="flex: 1; padding: 5px 8px; font-size: 12px; margin: 0;" onclick="editViolation(${v.id})">Edit Deskripsi</button>
                            <button class="btn-danger" style="flex: 1; padding: 5px 8px; font-size: 12px; margin: 0;" onclick="deleteViolation(${v.id})">Hapus Laporan</button>
                        </div>

          html += `<div style="border-bottom: 1px solid #ccc; padding: 10px 0;">
                            <p><strong>Penyewa:</strong> ${v.nama} (Unit ${v.unit})</p>
                            <p><strong>Deskripsi:</strong> <span id="desc-${v.id}">${v.description}</span></p>
                            ${v.photo_url ? `<img src="${v.photo_url}" style="display: block; margin: 10px auto; max-width: 100%; max-height: 300px; object-fit: contain; border-radius: 5px;" alt="Foto Pelanggaran">` : ''}
                            <p><small>Tanggal Laporan: ${new Date(v.created_at).toLocaleString('id-ID')}</small></p>
                            ${adminControls}
                        </div>;
        });
        container.innerHTML = html || '<p>Belum ada laporan.</p>';
      }
    }

    // Function to edit a violation report
    async function editViolation(violationId) {
      const currentDescription = document.getElementById(`desc-${violationId}`).innerText;
      const newDescription = prompt("Masukkan deskripsi baru:", currentDescription);
      if (newDescription === null || newDescription.trim() === '' || newDescription === currentDescription) return;

      const response = await fetch(`/api/dashboard/violations/${violationId}`, {
        method: 'PUT',
        headers: getAuthHeaders({ 'Content-Type': 'application/json' }),
        body: JSON.stringify({ description: newDescription })
      });
      const result = await response.json();
      showAlert(result.message, result.success ? 'success' : 'danger');
      if (result.success) loadDashboardViolations();
    }

    // Function to delete a violation report
    async function deleteViolation(violationId) {
      if (!confirm('Anda yakin ingin menghapus laporan pelanggaran ini? Tindakan ini tidak dapat dibatalkan.')) return;

      const response = await fetch(`/api/dashboard/violations/${violationId}`, {
        method: 'DELETE',
        headers: getAuthHeaders()
      });
      const result = await response.json();
      showAlert(result.message, result.success ? 'success' : 'danger');
      if (result.success) loadDashboardViolations();
    }


    async function loadRentalsForViolationDropdown() {
      const select = document.getElementById('violation-rental-id');
      const response = await fetch('/api/rentals', { headers: getAuthHeaders() });
      const result = await response.json();
      if (result.success) {
        select.innerHTML = '<option value="">Pilih Penyewa</option>';
        result.data.forEach(r => {
          select.innerHTML += `<option value="${r.id}">${r.nama} (Unit ${r.tower}-${r.lantai}-${r.unit})</option>`;
        });
      }
    }

    async function submitViolation() {
      if (!isUserAdmin(currentUser)) {
        showAlert('Akses ditolak. Hanya Pengelola yang dapat mengupload laporan.', 'danger');
        return;
      }
      const formData = new FormData();
      formData.append('rental_id', document.getElementById('violation-rental-id').value);
      formData.append('description', document.getElementById('violation-desc').value);
      const photoFile = document.getElementById('violation-photo-file').files[0];
      if (photoFile) {
        formData.append('photo', photoFile);
      }
      const response = await fetch('/api/dashboard/violations', { method: 'POST', headers: getAuthHeaders(), body: formData });
      const result = await response.json();
      showAlert(result.message, result.success ? 'success' : 'danger');
      if (result.success) {
        document.getElementById('violation-desc').value = '';
        document.getElementById('violation-photo-file').value = '';
        loadDashboardViolations();
      }
    }


    function renderPagination(totalItems, totalPages, page) {
      const containerBottom = document.getElementById('pagination-controls-bottom');

      if (totalPages <= 1) {
        if (containerBottom) containerBottom.innerHTML = '';
        return;
      }

      let paginationHTML = '';

      paginationHTML += `<button onclick="changePage(${page - 1})" ${page === 1 ? 'disabled' : ''}>&laquo; Prev</button>`;

      const maxPagesToShow = 5;
      let startPage = Math.max(1, page - Math.floor(maxPagesToShow / 2));
      let endPage = Math.min(totalPages, startPage + maxPagesToShow - 1);

      if (endPage - startPage + 1 < maxPagesToShow) {
        startPage = Math.max(1, endPage - maxPagesToShow + 1);
      }

      if (startPage > 1) {
        paginationHTML += `<button onclick="changePage(1)">1</button>`;
        if (startPage > 2) {
          paginationHTML += `<span class="page-ellipsis">...</span>`;
        }
      }

      for (let i = startPage; i <= endPage; i++) {
        paginationHTML += `<button onclick="changePage(${i})" class="${i === page ? 'active' : ''}">${i}</button>`;
      }

      if (endPage < totalPages) {
        if (endPage < totalPages - 1) {
          paginationHTML += `<span class="page-ellipsis">...</span>`;
        }
        paginationHTML += `<button onclick="changePage(${totalPages})">${totalPages}</button>`;
      }

      paginationHTML += `<button onclick="changePage(${page + 1})" ${page === totalPages ? 'disabled' : ''}>Next &raquo;</button>`;

      const startItem = (page - 1) * itemsPerPage + 1;
      const endItem = Math.min(page * itemsPerPage, totalItems);
      paginationHTML += `<span class="pagination-info">Menampilkan ${startItem}-${endItem} dari ${totalItems} data</span>`;

      if (containerBottom) containerBottom.innerHTML = paginationHTML;
    }


    function changePage(page) {
      // Tetapkan halaman baru dan muat ulang data dengan mode tampilan saat ini
      currentPage = page;
      loadRentalsData(currentViewMode, false);
    }

    async function updateRental(id) {
      const data = {
        jenis_sewa: document.getElementById(`edit_jenis_sewa_${id}`).value,
        metode_pembayaran: document.getElementById(`edit_metode_pembayaran_${id}`).value,
        waktu_checkout: document.getElementById(`edit_waktu_checkout_${id}`).value,
        komentar: Array.from(document.querySelectorAll(`.komentar-cb-${id}:checked`)).map(cb => cb.value)
      };
      const response = await fetch(`/api/rentals/${id}`, {
        method: 'PUT',
        headers: getAuthHeaders({ 'Content-Type': 'application/json' }),
        body: JSON.stringify(data)
      });
      const result = await response.json();
      showAlert(result.message, result.success ? 'success' : 'danger');

      // Muat ulang tampilan saat ini, reset ke halaman 1
      if (result.success) loadRentalsData(currentViewMode, true);

    }

    async function loadAgentsForManagement() {
      const select = document.getElementById('manage-agent-select');
      select.innerHTML = '<option>Memuat pengguna Pelaku Komersil...</option>';

      try {
        const response = await fetch('/api/agents', { headers: getAuthHeaders() });
        const result = await response.json();

        if (!result.success) {
          select.innerHTML = '<option value="">Gagal memuat data</option>';
          showAlert(result.message || 'Gagal memuat daftar pengguna Pelaku Komersil.', 'danger');
          return;
        }

        select.innerHTML = '<option value="">Pilih Pengguna Pelaku Komersil</option>';
        result.data.forEach(agent => {
          const optionValue = `id:${agent.id}`;
          const optionLabel = agent.nama
            ? `${agent.nama} (${agent.email || 'tanpa email'})`
            : (agent.email || `Pengguna #${agent.id}`);
          select.innerHTML += `<option value="${optionValue}">${optionLabel}</option>`;
        });
      } catch (error) {
        console.error('Error loading Pelaku Komersil users:', error);
        select.innerHTML = '<option value="">Gagal memuat data</option>';
        showAlert('Gagal terhubung ke server untuk memuat pengguna Pelaku Komersil.', 'danger');
      }
    }

    async function loadUnitsForAgent() {
      const agentKey = document.getElementById('manage-agent-select').value;
      const container = document.getElementById('unit-list-container');
      if (!agentKey) { container.innerHTML = ''; return; }
      container.innerHTML = 'Memuat unit...';

      const response = await fetch(`/api/units/${encodeURIComponent(agentKey)}`, { headers: getAuthHeaders() });
      const result = await response.json();
      if (result.success) {
        let html = '<ul>';
        result.data.forEach(unit => {
          html += `<li>${unit.unit_number} <button onclick="deleteUnit('${unit.id}')" style="background: #dc3545; font-size: 12px; padding: 2px 5px;">Hapus</button></li>`;
        });
        container.innerHTML = (result.data.length > 0 ? html + '</ul>' : '<p>Pengguna ini belum memiliki unit.</p>');
      }
    }

    async function addUnit() {
      if (!isUserAdmin(currentUser)) {
        showAlert('Akses ditolak. Hanya Pengelola yang dapat menambah unit.', 'danger');
        return;
      }
      const data = {
        agent_key: document.getElementById('manage-agent-select').value,
        unit_tower: document.getElementById('new-unit-tower').value,
        unit_lantai: document.getElementById('new-unit-lantai').value,
        unit_nomor: document.getElementById('new-unit-nomor').value,
      };
      if (!data.agent_key || !data.unit_tower || !data.unit_lantai || !data.unit_nomor) {
        showAlert('Semua field unit harus diisi.', 'danger'); return;
      }
      const response = await fetch('/api/units', {
        method: 'POST',
        headers: getAuthHeaders({ 'Content-Type': 'application/json' }),
        body: JSON.stringify(data)
      });
      const result = await response.json();
      showAlert(result.message, result.success ? 'success' : 'danger');
      if (result.success) loadUnitsForAgent();
    }

    async function deleteUnit(unitId) {
      console.log('Menghapus unit dengan ID:', unitId);
      if (!confirm('Anda yakin ingin menghapus unit ini?')) return;
      const response = await fetch(`/api/units/${unitId}`, { method: 'DELETE', headers: getAuthHeaders() });
      const result = await response.json();
      showAlert(result.message, result.success ? 'success' : 'danger');
      if (result.success) loadUnitsForAgent();
    }

    async function loadMyUnitsForInput() {
      const select = document.getElementById('unit-select');
      select.innerHTML = '<option value="">Memuat unit...</option>';
      const isAdmin = isUserAdmin(currentUser);
      const url = isAdmin ? '/api/all-units' : '/api/my-units';

      try {
        const response = await fetch(url, { headers: getAuthHeaders() });
        const result = await response.json();
        if (result.success) {
          let options = '<option value="">-- Pilih Unit --</option>';
          result.data.forEach(u => {
            let displayText = isAdmin
              ? `Tower ${u.unit_tower} - Lt ${u.unit_lantai} - No ${u.unit_nomor} (Agen: ${u.agent_email})`
              : `Tower ${u.unit_tower} - Lt ${u.unit_lantai} - No ${u.unit_nomor}`;

            let isDisabled = u.is_occupied;
            if (isDisabled) {
              displayText += ' (Sedang Disewa)';
            }

            options += `<option value="${u.unit_number}" ${isDisabled ? 'disabled' : ''}>${displayText}</option>`;
          });
          select.innerHTML = options;
        } else {
          select.innerHTML = `<option value="">Gagal memuat unit: ${result.message}</option>`;
        }
      } catch (error) {
        select.innerHTML = '<option value="">Error saat memuat data</option>';
      }
    }

    function populateUnitFields() {
      const selectedUnit = document.getElementById('unit-select').value;
      if (selectedUnit) {
        const [tower, lantai, nomor] = selectedUnit.split('-');
        document.getElementById('tower').value = tower;
        document.getElementById('lantai').value = lantai;
        document.getElementById('unit').value = nomor;
      }
    }

    function showAlert(message, type) {
      const container = document.getElementById('alert-container');
      const alertDiv = document.createElement('div');
      alertDiv.className = `alert alert-${type}`;
      alertDiv.textContent = message;
      container.appendChild(alertDiv);
      setTimeout(() => alertDiv.remove(), 5000);
    }

    function showPredictionNotification(message) {
      const container = document.getElementById('alert-container');
      const alertDiv = document.createElement('div');
      alertDiv.className = 'alert';
      alertDiv.style.backgroundColor = '#fff3cd';
      alertDiv.style.borderColor = '#ffe69c';
      alertDiv.style.color = '#664d03';
      alertDiv.style.boxShadow = '0 4px 14px rgba(0, 0, 0, 0.08)';
      alertDiv.style.fontWeight = '600';
      alertDiv.textContent = message;
      container.appendChild(alertDiv);
      setTimeout(() => alertDiv.remove(), 5000);
    }
  </script>
</body>

</html>
\end{lstlisting}





\begin{lstlisting}[language=JavaScript, caption={daftar.html}]
<!DOCTYPE html>
<html lang="id">
<head>
  <meta charset="UTF-8">
  <meta name="viewport" content="width=device-width, initial-scale=1.0">
  <title>Daftar - Data Capture</title>
  <link href="https://fonts.googleapis.com/css2?family=Montserrat:wght@400;500;600;700&display=swap" rel="stylesheet">
  <style>
    :root {
      --theme: #399051;
      --themeHover: #2e8546;
    }
    * {
      font-family: "Montserrat", sans-serif !important;
      box-sizing: border-box;
    }
    body {
      margin: 0;
      padding: 0;
      background-color: #f3f7fb;
      display: flex;
      justify-content: center;
      align-items: center;
      height: 100vh;
    }
    .login-container {
      background: white;
      padding: 30px;
      border-radius: 8px;
      box-shadow: 0px 4px 25px 0px rgba(0, 0, 0, 0.06);
      width: 100%;
      max-width: 400px;
    }
    .login-container h2 {
      text-align: center;
      color: #333;
      margin-bottom: 20px;
    }
    .form-group {
      margin-bottom: 15px;
    }
    .form-group label {
      display: block;
      margin-bottom: 5px;
      font-size: 14px;
      color: #555;
      font-weight: 500;
    }
    .form-group input {
      width: 100%;
      padding: 10px;
      border: 1px solid #ccc;
      border-radius: 4px;
      font-size: 14px;
    }
    .btn-success {
      width: 100%;
      background-color: var(--theme);
      color: white;
      border: none;
      padding: 12px;
      border-radius: 4px;
      cursor: pointer;
      font-size: 16px;
      font-weight: 600;
      margin-top: 10px;
      transition: background-color 0.2s;
    }
    .btn-success:hover {
      background-color: var(--themeHover);
    }
    .alert {
      padding: 10px;
      margin-bottom: 15px;
      border-radius: 4px;
      font-size: 14px;
      display: none;
    }
    .alert-danger {
      background-color: #f8d7da;
      color: #721c24;
    }
    .alert-success {
      background-color: #d4edda;
      color: #155724;
    }
    .text-center {
      text-align: center;
      margin-top: 15px;
      font-size: 14px;
    }
    .text-center a {
      color: var(--theme);
      text-decoration: none;
      font-weight: 500;
    }
  </style>
</head>
<body>
  <div class="login-container">
    <h2>Daftar Akun</h2>
    <div id="alert-message" class="alert"></div>
    <form id="register-form">
      <div class="form-group">
        <label for="name">Nama Lengkap</label>
        <input type="text" id="name" required>
      </div>
      <div class="form-group">
        <label for="email">Email</label>
        <input type="email" id="email" required>
      </div>
      <div class="form-group">
        <label for="password">Password</label>
        <input type="password" id="password" required>
      </div>
      <div class="form-group">
        <label for="nib">NIB (13 Digit)</label>
        <input type="text" id="nib" minlength="13" maxlength="13" required>
      </div>
      <div class="form-group">
        <label for="role">Role</label>
        <select id="role" style="width: 100%; padding: 10px; border: 1px solid #ccc; border-radius: 4px; font-size: 14px;" required>
          <!-- <option value="">Pilih Role</option>
          <option value="2">Pengelola</option> -->
          <option value="4">Pelaku Komersil</option>
        </select>
      </div>
      <button type="submit" class="btn-success">Daftar</button>
    </form>
    <div class="text-center">
      Sudah punya akun? <a href="/login.html">Login disini</a>
    </div>
  </div>

  <script>
    document.getElementById('register-form').addEventListener('submit', async function(e) {
      e.preventDefault();
      const name = document.getElementById('name').value;
      const email = document.getElementById('email').value;
      const password = document.getElementById('password').value;
      const nib = document.getElementById('nib').value;
      const role = document.getElementById('role').value;
      const alertDiv = document.getElementById('alert-message');
      
      try {
        const response = await fetch('/api/register', {
          method: 'POST',
          headers: {
            'Content-Type': 'application/json'
          },
          body: JSON.stringify({ name, email, password, nib, role })
        });
        const data = await response.json();
        
        if (data.success) {
          alertDiv.innerText = data.message;
          alertDiv.className = 'alert alert-success';
          alertDiv.style.display = 'block';
          setTimeout(() => {
            window.location.href = '/login.html';
          }, 2000);
        } else {
          alertDiv.innerText = data.message || 'Registrasi gagal';
          alertDiv.className = 'alert alert-danger';
          alertDiv.style.display = 'block';
        }
      } catch (err) {
        alertDiv.innerText = 'Terjadi kesalahan pada server';
        alertDiv.className = 'alert alert-danger';
        alertDiv.style.display = 'block';
      }
    });
  </script>
</body>
</html>
\end{lstlisting}


\begin{lstlisting}[language=JavaScript, caption={login.html}]
<!DOCTYPE html>
<html lang="id">
<head>
  <meta charset="UTF-8">
  <meta name="viewport" content="width=device-width, initial-scale=1.0">
  <title>Login - Data Capture</title>
  <link href="https://fonts.googleapis.com/css2?family=Montserrat:wght@400;500;600;700&display=swap" rel="stylesheet">
  <style>
    :root {
      --theme: #399051;
      --themeHover: #2e8546;
    }
    * {
      font-family: "Montserrat", sans-serif !important;
      box-sizing: border-box;
    }
    body {
      margin: 0;
      padding: 0;
      background-color: #f3f7fb;
      display: flex;
      justify-content: center;
      align-items: center;
      height: 100vh;
    }
    .login-container {
      background: white;
      padding: 30px;
      border-radius: 8px;
      box-shadow: 0px 4px 25px 0px rgba(0, 0, 0, 0.06);
      width: 100%;
      max-width: 400px;
    }
    .login-container h2 {
      text-align: center;
      color: #333;
      margin-bottom: 20px;
    }
    .form-group {
      margin-bottom: 15px;
    }
    .form-group label {
      display: block;
      margin-bottom: 5px;
      font-size: 14px;
      color: #555;
      font-weight: 500;
    }
    .form-group input {
      width: 100%;
      padding: 10px;
      border: 1px solid #ccc;
      border-radius: 4px;
      font-size: 14px;
    }
    .btn-success {
      width: 100%;
      background-color: var(--theme);
      color: white;
      border: none;
      padding: 12px;
      border-radius: 4px;
      cursor: pointer;
      font-size: 16px;
      font-weight: 600;
      margin-top: 10px;
      transition: background-color 0.2s;
    }
    .btn-success:hover {
      background-color: var(--themeHover);
    }
    .alert {
      padding: 10px;
      margin-bottom: 15px;
      border-radius: 4px;
      font-size: 14px;
      display: none;
    }
    .alert-danger {
      background-color: #f8d7da;
      color: #721c24;
    }
    .text-center {
      text-align: center;
      margin-top: 15px;
      font-size: 14px;
    }
    .text-center a {
      color: var(--theme);
      text-decoration: none;
      font-weight: 500;
    }
  </style>
</head>
<body>
  <div class="login-container">
    <h2>Login</h2>
    <div id="error-message" class="alert alert-danger"></div>
    <form id="login-form">
      <div class="form-group">
        <label for="email">Email</label>
        <input type="email" id="email" required>
      </div>
      <div class="form-group">
        <label for="password">Password</label>
        <input type="password" id="password" required>
      </div>
      <button type="submit" class="btn-success">Masuk</button>
    </form>
    <div class="text-center">
      Belum punya akun? <a href="/daftar.html">Daftar disini</a>
    </div>
  </div>

  <script>
    document.getElementById('login-form').addEventListener('submit', async function(e) {
      e.preventDefault();
      const email = document.getElementById('email').value;
      const password = document.getElementById('password').value;
      const errorDiv = document.getElementById('error-message');
      
      try {
        // Correct implementation ?
        const response = await fetch('/api/login', {
            method: 'POST',
            headers: { 'Content-Type': 'application/json' },
            body: JSON.stringify({ 
                email: email, 
                password: password 
            })
        });
        const data = await response.json();
        
        if (data.success) {
          localStorage.setItem('userSession', JSON.stringify({ token: data.token, object: data.user }));
          window.location.href = '/';
        } else {
          errorDiv.innerText = data.message || 'Login gagal';
          errorDiv.style.display = 'block';
        }
      } catch (err) {
        errorDiv.innerText = 'Terjadi kesalahan pada server';
        errorDiv.style.display = 'block';
      }
    });
  </script>
</body>
</html>
\end{lstlisting}



\begin{lstlisting}[language=JavaScript, caption={app.js}]
// // package.json dependencies needed:
// // npm install express mysql2 dotenv path body-parser cors multer json2csv jsonwebtoken

const express = require('express');
const mysql = require('mysql2/promise');
const path = require('path');
const bodyParser = require('body-parser');
const jwt = require('jsonwebtoken');
const bcrypt = require('bcryptjs');
const cors = require('cors');
const multer = require('multer');
const { spawnSync } = require('child_process');
const { Parser } = require('json2csv');
require('dotenv').config({
  path: path.resolve(__dirname, '.env'),
  override: true
});

const app = express();
const port = process.env.PORT || 3000;
const JWT_SECRET = process.env.JWT_SECRET || 'jwt-secret-key-jarrdin';
const SRUSUN_JWT_SECRET = process.env.SRUSUN_JWT_SECRET || 'jarrdin-cihampelas';
const MODEL_PATH = path.resolve(__dirname, 'model_xgboost.pkl');
const PYTHON_EXECUTABLE = process.env.PYTHON_EXECUTABLE || path.resolve(__dirname, '.venv', 'Scripts', 'python.exe');

// --- File Upload Configuration ---
const storage = multer.diskStorage({
  destination: function (req, file, cb) {
    cb(null, 'public/uploads/');
  },
  filename: function (req, file, cb) {
    const uniqueSuffix = Date.now() + '-' + Math.round(Math.random() * 1E9);
    cb(null, file.fieldname + '-' + uniqueSuffix + path.extname(file.originalname));
  }
});
const upload = multer({ storage: storage });

// Middleware
const allowedOrigins = [
  'https://web.srusun.id',
  'https://webmember.srusun.id',
  'https://member.thejarrdin.com',
  'https://apimember.thejarrdin.com',
  'https://visitor.thejarrdin.com',
  'http://localhost:5173',
  'http://localhost:3000',
  'http://localhost:3001',
  'http://127.0.0.1:5173'
];

const corsOptions = {
  origin: (origin, callback) => {
    if (!origin || allowedOrigins.includes(origin)) {
      return callback(null, true);
    }
    // Don't throw an error, just block the request nicely
    return callback(null, false);
  },
  credentials: true,
  methods: ['GET', 'POST', 'PUT', 'PATCH', 'DELETE', 'OPTIONS'],
  allowedHeaders: ['Content-Type', 'Authorization']
};

app.use(cors(corsOptions));
app.options('*', cors(corsOptions));
app.use(bodyParser.json());
app.use(bodyParser.urlencoded({ extended: true }));
app.use(express.static('public'));
app.use('/uploads', express.static('public/uploads'));

// Remote SRusun DB - used for: rentals, rental_logs, units, violations
const pool = mysql.createPool({
  host: process.env.SRUSUN_MYSQL_HOST,
  user: process.env.SRUSUN_MYSQL_USER,
  password: process.env.SRUSUN_MYSQL_PASSWORD,
  database: process.env.SRUSUN_MYSQL_DATABASE,
  waitForConnections: true,
  connectionLimit: 10,
  queueLimit: 0,
  enableKeepAlive: true,
  keepAliveInitialDelay: 10000
});

// --- Security Middleware ---
const checkAuth = (req, res, next) => {
  const authHeader = req.headers['authorization'];
  const token = (authHeader && authHeader.split(' ')[1]) || req.query.token;
  if (!token) {
    return res.status(401).json({ success: false, message: "Akses ditolak. Silakan login." });
  }

  let decoded;
  try {
    decoded = jwt.verify(token, JWT_SECRET);
  } catch (_) {
    try {
      decoded = jwt.verify(token, SRUSUN_JWT_SECRET);
    } catch {
      return res.status(401).json({ success: false, message: "Token tidak valid. Silakan login ulang." });
    }
  }

  // Normalize: SRusun JWT has Role (array of objects); Data Capture JWT has roles (array of strings)
  // Use Array.isArray(decoded.Role) as the primary SRusun detector - more robust than checking UserID
  if (Array.isArray(decoded.Role)) {
    req.user = {
      id: decoded.UserID ?? decoded.userId,
      nama: decoded.Nama,
      email: decoded.Email,
      no_telp: decoded.NoTelp,
      roles: decoded.Role.map(r => (typeof r === 'object' ? r.Nama : r)),
      fitur: decoded.Fitur || [],
      _source: 'srusun'
    };
  } else {
    req.user = decoded;
  }
  return next();
};

// SRusun roles: 'Pengelola' has admin access, 'Pelaku Komersil' has regular access
const ADMIN_ROLES = ['Pengelola'];
const isAdmin = (user) => {
  if (!user) return false;
  // Handle normalized format: roles is array of strings (set by checkAuth normalization)
  if (user.roles && user.roles.some(r => ADMIN_ROLES.includes(r))) return true;
  // Handle raw SRusun format: Role is array of objects/strings (fallback if normalization didn't run)
  if (Array.isArray(user.Role) && user.Role.some(r => ADMIN_ROLES.includes(typeof r === 'object' ? r.Nama : r))) return true;
  return false;
};
const checkAdmin = (req, res, next) => {
  checkAuth(req, res, () => {
    if (!isAdmin(req.user)) {
      return res.status(403).json({ success: false, message: "Akses ditolak. Fitur ini hanya untuk admin." });
    }
    next();
  });
};

// ================== AUTHENTICATION ==================
app.post('/api/login', async (req, res) => {
  // 1. Safely extract and trim the email to prevent trailing space errors
  const email = req.body.email ? req.body.email.trim() : undefined;
  const password = req.body.password;

  // 2. Validate that the frontend actually sent the email!
  if (!email || !password) {
    return res.status(400).json({ 
      success: false, 
      message: 'Email atau password kosong. Cek payload dari login.html!' 
    });
  }

  try {
    // Arahkan query ke tabel Pengguna, bukan users
    const [rows] = await pool.query('SELECT * FROM Pengguna WHERE email = ?', [email]);
    
    if (rows.length === 0) {
      return res.status(401).json({ success: false, message: 'Email tidak ditemukan.' });
    }
    
    const user = rows[0];
    
    // Verifikasi Password 
    const isMatch = await bcrypt.compare(password, user.password); 
    if (!isMatch) {
      return res.status(401).json({ success: false, message: 'Password salah.' });
    }

    // Ambil role user dari tabel Role agar middleware checkAdmin () berfungsi
    const [roleRows] = await pool.query(`
      SELECT r.nama as role_name
      FROM pengguna_role pr
      JOIN Role r ON r.role_id = pr.role_id
      WHERE pr.pengguna_id = ?
    `, [user.pengguna_id]);

    const roles = roleRows.map(r => r.role_name);

    // Susun payload menggunakan ID dari Pengguna
    const payload = { 
      id: user.pengguna_id, 
      email: user.email, 
      nama: user.nama, 
      roles: roles 
    };
    
    const token = jwt.sign(payload, JWT_SECRET, { expiresIn: '12h' });
    res.json({ success: true, token, user: payload });
    
  } catch (error) {
    console.error('Login Error:', error);
    res.status(500).json({ success: false, message: error });
  }
});

app.post('/api/register', async (req, res) => {
  const { name, email, password, nib, role } = req.body;
  
  if (!email || !password || !nib || !name || !role) {
    return res.status(400).json({ success: false, message: 'Semua kolom harus diisi.' });
  }
  if (nib.length !== 13) {
    return res.status(400).json({ success: false, message: 'NIB harus 13 digit.' });
  }
  
  const conn = await pool.getConnection();
  try {
    // Check existing email purely in Pengguna table
    const [existing] = await conn.query('SELECT * FROM Pengguna WHERE email = ?', [email]);
    if (existing.length > 0) {
      conn.release();
      return res.status(400).json({ success: false, message: 'Email sudah terdaftar.' });
    }
    
    const hashedPassword = await bcrypt.hash(password, 10);
    
    await conn.beginTransaction();

    // 1. Manually find the current highest pengguna_id in the table
    const [idRows] = await conn.query('SELECT MAX(pengguna_id) as max_id FROM Pengguna');
    
    // 2. Add 1 to create the next ID (if the table is empty, start at 1)
    const nextId = (idRows[0].max_id || 0) + 1;
    
    // 3. Inject nextId directly into the INSERT query for Pengguna
    await conn.query(
      'INSERT INTO Pengguna (pengguna_id, kode_user, nama, no_unit, email, password, nib, status, is_ketentuan) VALUES (?, ?, ?, ?, ?, ?, ?, 1, 1)', 
      [nextId, name, name, '', email, hashedPassword, nib]
    );
    
    // 4. Link the new user to their role in pengguna_role
    await conn.query('INSERT INTO pengguna_role (pengguna_id, role_id) VALUES (?, ?)', [nextId, role]);
    
    // (The legacy users table insertion has been completely removed)

    await conn.commit();
    conn.release();
    res.json({ success: true, message: 'Registrasi berhasil, silakan login.' });
    
  } catch (error) {
    if (conn) await conn.rollback();
    if (conn) conn.release();
    console.error('Register Error:', error);
    res.status(500).json({ success: false, message: 'Server error.' });
  }
});

app.post('/api/logout', (req, res) => {
  res.json({ success: true });
});

app.get('/api/user', checkAuth, (req, res) => {
  res.json({ user: req.user });
});

// Alias for session verification used by the React client
app.get('/api/user/session', checkAuth, (req, res) => {
  res.json({ user: req.user });
});


// ================== RENTALS DATA ==================
app.get('/api/rentals', checkAuth, async (req, res) => {
  try {
    // PAGINATION & FILTER LOGIC
    const page = parseInt(req.query.page) || 1;
    const limit = parseInt(req.query.limit) || 10;
    const hasCommentFilter = req.query.commentFilter === 'true';
    const searchBox = req.query.search || '';
    const offset = (page - 1) * limit;

    const user = req.user;
    let whereClauses = [];
    let params = [];

    // Filter berdasarkan Role
    if (!isAdmin(user)) {
      whereClauses.push('r.user_pengguna_id = ?');
      params.push(user.id);
    }

    // Filter berdasarkan Komentar (jika ada)
    if (hasCommentFilter) {
      whereClauses.push("(r.komentar IS NOT NULL AND r.komentar != '[]' AND r.komentar != '')");
    }

    // Filter berdasarkan pencarian
    if (searchBox) {
      const searchPattern = `%${searchBox}%`;
      whereClauses.push("(r.nama LIKE ? OR r.nik LIKE ? OR r.unit_number LIKE ?)");
      params.push(searchPattern, searchPattern, searchPattern);
    } // FIXED: Added missing closing bracket here

    // FIXED: Defined whereSql before the queries
    const whereSql = whereClauses.length > 0 ? `WHERE ${whereClauses.join(' AND ')}` : '';

    // --- Query 1: Menghitung Total Item ---
    const countQuery = `SELECT COUNT(*) as totalItems FROM rentals r ${whereSql}`;
    const [countRows] = await pool.query(countQuery, params);
    const totalItems = countRows[0].totalItems;
    const totalPages = Math.ceil(totalItems / limit);

    // --- Query 2: Mengambil Data untuk Halaman Ini ---
    const logSubQuery = `(SELECT GROUP_CONCAT(
            JSON_OBJECT(
                'action', l.action, 'field_changed', l.field_changed, 'old_value', l.old_value,
                'new_value', l.new_value, 'email', lp.email, 'timestamp', l.timestamp
            ) SEPARATOR '|||'
        ) FROM rental_logs l LEFT JOIN Pengguna lp ON l.pengguna_id = lp.pengguna_id WHERE l.rental_id = r.id) as logs_json`;

    const dataQuery = `
            SELECT r.*, u.tower, u.lantai, u.unit, p.email AS user_email, e.email AS diedit_oleh, r.user_pengguna_id as user_pengguna_id, ${logSubQuery}
            FROM rentals r
            LEFT JOIN units u ON r.unit_number = u.unit_number
            LEFT JOIN Pengguna p ON r.user_pengguna_id = p.pengguna_id
            LEFT JOIN Pengguna e ON r.editor_pengguna_id = e.pengguna_id
            ${whereSql}
            ORDER BY r.created_at DESC
            LIMIT ? OFFSET ?
        `;

    // Tambahkan parameter limit dan offset ke params
    const dataParams = [...params, limit, offset];
    await pool.query('SET SESSION group_concat_max_len = 1000000');
    const [rows] = await pool.query(dataQuery, dataParams);
    
    rows.forEach(r => {
      r.status_pasutri = r.status_pasutri === 'Ya' ? 'Ya' : 'Belum Menikah';
      try {
        r.komentar = JSON.parse(r.komentar || '[]');
      } catch (e) {
        r.komentar = r.komentar ? [r.komentar] : [];
      }
      if (r.logs_json) {
        r.logs = r.logs_json.split('|||').map(logStr => {
          try {
            return JSON.parse(logStr);
          } catch (e) {
            console.log(r)
            console.log("Offending log string:", logStr);
            console.log("Offending log string:", r.logs_json);
            return { action: "invalid_log", error: "Invalid JSON" };
          }
        }).sort((a, b) => {
          if (!a.timestamp || !b.timestamp) return 0;
          return new Date(b.timestamp) - new Date(a.timestamp);
        });
      } else {
        r.logs = [];
      }
      delete r.logs_json;
    });

    // Kirim Respon Terstruktur Baru 
    res.json({
      success: true,
      data: rows,
      pagination: {
        page: page,
        limit: limit,
        totalItems: totalItems,
        totalPages: totalPages
      }
    });

  } catch (error) {
    console.error("Error fetching rentals with logs:", error);
    res.status(500).json({ success: false, message: 'Gagal mengambil data penyewa.' });
  }
});

app.put('/api/rentals/:id', checkAuth, async (req, res) => {
  const { id } = req.params;
  const { jenis_sewa, metode_pembayaran, waktu_checkout, komentar, metode_lain } = req.body;

  const conn = await pool.getConnection();
  try {
    await conn.beginTransaction();

    const [oldDataRows] = await conn.query(`
            SELECT r.*, u.* FROM rentals r LEFT JOIN units u ON r.unit_number = u.unit_number WHERE r.id = ?
            `, [id]);
    if (oldDataRows.length === 0) {
      await conn.rollback(); conn.release();
      return res.status(404).json({ success: false, message: 'Data rental tidak ditemukan.' });
    }
    const oldData = oldDataRows[0];

    const newData = {
      jenis_sewa,
      metode_pembayaran,
      metode_lain: metode_lain || oldData.metode_lain,
      waktu_checkout: waktu_checkout || null,
      komentar: JSON.stringify(komentar || []),
    };
    const changedFields = {};
    const logEntries = [];

    Object.keys(newData).forEach(key => {
      let oldValue = oldData[key];
      if (key === 'komentar') oldValue = oldValue || '[]';

      if (String(oldValue) !== String(newData[key])) {
        changedFields[key] = newData[key];
        logEntries.push(['update', key, String(oldValue), String(newData[key]), req.user.id, id]);
      }
    });

    if (Object.keys(changedFields).length > 0) {
      changedFields.editor_pengguna_id = req.user.id;
      const updateQuery = 'UPDATE rentals SET ' + Object.keys(changedFields).map(key => `${key} = ?`).join(', ') + ' WHERE id = ?';
      const updateParams = [...Object.values(changedFields), id];
      await conn.query(updateQuery, updateParams);

      if (logEntries.length > 0) {
        const logQuery = 'INSERT INTO rental_logs (action, field_changed, old_value, new_value, pengguna_id, rental_id) VALUES ?';
        await conn.query(logQuery, [logEntries]);
      }
    }

    await conn.commit();
    conn.release();
    res.json({ success: true, message: 'Data berhasil diperbarui.' });

  } catch (error) {
    if (conn) await conn.rollback();
    if (conn) conn.release();
    console.error("Error updating rental:", error);
    res.status(500).json({ success: false, message: 'Gagal memperbarui data.' });
  }
});

app.get('/api/rentals/duplicates', checkAdmin, async (req, res) => {
  try {

    const page = parseInt(req.query.page) || 1;
    const limit = parseInt(req.query.limit) || 10;
    const offset = (page - 1) * limit;

    const innerQuery = `(SELECT nik FROM rentals WHERE nik IS NOT NULL AND nik != '' GROUP BY nik HAVING COUNT(*) > 1) r2`;

    // --- Query 1: Menghitung Total Item ---
    const countQuery = `SELECT COUNT(DISTINCT r1.id) as totalItems FROM rentals r1 INNER JOIN ${innerQuery} ON r1.nik = r2.nik`;
    const [countRows] = await pool.query(countQuery);
    const totalItems = countRows[0].totalItems;
    const totalPages = Math.ceil(totalItems / limit);

    // --- Query 2: Mengambil Data ---
    const dataQuery = `
            SELECT r1.* FROM rentals r1
            INNER JOIN ${innerQuery} ON r1.nik = r2.nik
            ORDER BY r1.nik, r1.created_at DESC
            LIMIT ? OFFSET ?
        `;
    const [rows] = await pool.query(dataQuery, [limit, offset]);


    rows.forEach(r => r.komentar = JSON.parse(r.komentar || '[]'));

    // Kirim Respon Terstruktur Baru 
    res.json({
      success: true,
      data: rows,
      pagination: {
        page: page,
        limit: limit,
        totalItems: totalItems,
        totalPages: totalPages
      }
    });
  } catch (error) {
    res.status(500).json({ success: false, message: 'Gagal cek data duplikat.' });
  }
});


// ================== CSV DOWNLOAD ENDPOINTS ==================

// --- 1. DOWNLOAD RENTAL DATA (All Roles) ---
app.get('/api/rentals/csv', checkAuth, async (req, res) => {
  try {
    const user = req.user;
    let query = 'SELECT r.*, u.tower, u.lantai, u.unit, p.email AS user_email, e.email AS diedit_oleh FROM rentals r LEFT JOIN units u ON r.unit_number = u.unit_number LEFT JOIN Pengguna p ON r.user_pengguna_id = p.pengguna_id LEFT JOIN Pengguna e ON r.editor_pengguna_id = e.pengguna_id';
    const params = [];

    if (!isAdmin(user)) {
      query += ' WHERE r.user_pengguna_id = ?';
      params.push(user.id);
    }
    query += ' ORDER BY r.created_at DESC';
    const [rows] = await pool.query(query, params);

    const fields = ['id', 'nama', 'nik', 'status_pasutri', 'status_kewarganegaraan', 'tower', 'lantai', 'unit', 'jenis_sewa', 'metode_pembayaran', 'metode_lain', 'tanggal_checkin', 'waktu_checkin', 'tanggal_checkout', 'waktu_checkout', 'lama_menginap', 'user_email', 'diedit_oleh', 'created_at'];
    const csv = new Parser({ fields }).parse(rows);

    res.header('Content-Type', 'text/csv');
    res.attachment('data-penyewa.csv');
    res.send(csv);
  } catch (error) {
    res.status(500).send("Gagal membuat file CSV.");
  }
});

// --- 2. DOWNLOAD ALL LOG DATA (All Roles) ---
app.get('/api/logs/csv', checkAuth, async (req, res) => {
  try {
    const user = req.user;
    let query = `
            SELECT 
                l.id, l.timestamp, l.pengguna_id, p.email, r.nama as nama_penyewa, CONCAT(u.tower, '-', u.lantai, '-', u.unit) as unit,
                l.action, l.field_changed, l.old_value, l.new_value
            FROM rental_logs l
            LEFT JOIN rentals r ON l.rental_id = r.id
            LEFT JOIN units u ON r.unit_number = u.unit_number
            LEFT JOIN Pengguna p ON l.pengguna_id = p.pengguna_id
        `;
    const params = [];

    if (!isAdmin(user)) {
      query += ' WHERE r.user_pengguna_id = ?';
      params.push(user.id);
    }

    query += ' ORDER BY l.timestamp DESC';

    const [logs] = await pool.query(query, params);

    const fields = ['id', 'timestamp', 'email', 'nama_penyewa', 'unit', 'action', 'field_changed', 'old_value', 'new_value'];
    const csv = new Parser({ fields }).parse(logs);

    res.header('Content-Type', 'text/csv');
    res.attachment('semua-log-perubahan.csv');
    res.send(csv);
  } catch (error) {
    res.status(500).send("Gagal membuat file CSV.");
  }
});

// --- 3. DOWNLOAD SINGLE TENANT LOG (All Roles) ---
app.get('/api/rentals/:id/log/csv', checkAuth, async (req, res) => {
  try {
    const { id } = req.params;

    const [rentalInfo] = await pool.query('SELECT r.nama, CONCAT(u.tower, \'-\', u.lantai, \'-\', u.unit) as unit FROM rentals r LEFT JOIN units u ON r.unit_number = u.unit_number WHERE r.id = ?', [id]);
    if (rentalInfo.length === 0) return res.status(404).send('Data penyewa tidak ditemukan.');

    const { nama, unit } = rentalInfo[0];
    const safeFilename = `log-${nama.replace(/[^a-zA-Z0-9]/g, '_')}-unit_${unit}.csv`;

    const [logs] = await pool.query(
      'SELECT l.*, p.email FROM rental_logs l LEFT JOIN Pengguna p ON l.pengguna_id = p.pengguna_id WHERE l.rental_id = ? ORDER BY l.timestamp DESC', 
      [id]
    );

    const fields = ['timestamp', 'email', 'action', 'field_changed', 'old_value', 'new_value'];
    const csv = new Parser({ fields }).parse(logs);

    res.header('Content-Type', 'text/csv');
    res.attachment(safeFilename);
    res.send(csv);
  } catch (error) {
    res.status(500).send("Gagal membuat file CSV.");
  }
});


// ================== DASHBOARD (NEW & UPDATED) ==================

// --- 1. RENTALS OVER TIME (FOR LINE CHART) ---
app.get('/api/dashboard/rentals-over-time', checkAuth, async (req, res) => {
  const { filter } = req.query;
  const user = req.user;

  let groupBy = "DATE_FORMAT(tanggal_checkin, '%Y-%m-%d')"; // Daily for weekly
  let dateFilter = 'WHERE tanggal_checkin >= DATE_SUB(NOW(), INTERVAL 7 DAY)';

  if (filter === 'monthly') {
    dateFilter = 'WHERE tanggal_checkin >= DATE_SUB(NOW(), INTERVAL 1 MONTH)';
  } else if (filter === 'all') {
    groupBy = "DATE_FORMAT(tanggal_checkin, '%Y-%m')"; // Monthly for all-time
    dateFilter = '';
  }

  let query = `
        SELECT ${groupBy} as period, COUNT(*) as count 
        FROM rentals 
        ${dateFilter}
    `;
  const params = [];

  if (!isAdmin(user)) {
    query += ` ${dateFilter ? 'AND' : 'WHERE'} user_pengguna_id = ?`;
    params.push(user.id);
  }

  query += ` GROUP BY period ORDER BY period ASC`;

  try {
    const [results] = await pool.query(query, params);
    const labels = results.map(r => r.period);
    const data = results.map(r => r.count);
    res.json({ success: true, labels, data });
  } catch (error) {
    console.error("Stats Error:", error);
    res.status(500).json({ success: false, message: 'Gagal memuat statistik rentang waktu.' });
  }
});

// --- 2. POPULAR UNITS (FOR HISTOGRAM) ---
app.get('/api/dashboard/popular-units', checkAuth, async (req, res) => {
  const { filter } = req.query;
  const user = req.user;

  let dateFilter = '';
  if (filter === 'weekly') {
    dateFilter = 'WHERE created_at >= DATE_SUB(NOW(), INTERVAL 7 DAY)';
  } else if (filter === 'monthly') {
    dateFilter = 'WHERE created_at >= DATE_SUB(NOW(), INTERVAL 1 MONTH)';
  }

  // userEmail will only exist if we map it properly, but here we group by user_pengguna_id.
  let query = `
        SELECT 
            CONCAT(u.tower, '-', u.lantai, '-', u.unit) as unit_full, 
            r.user_pengguna_id as user_pengguna_id,
            COUNT(*) as rental_count 
        FROM rentals r
        LEFT JOIN units u ON r.unit_number = u.unit_number
        ${dateFilter}
    `;
  const params = [];

  if (!isAdmin(user)) {
    query += ` ${dateFilter ? 'AND' : 'WHERE'} r.user_pengguna_id = ?`;
    params.push(user.id);
  }

  query += ` GROUP BY unit_full, r.user_pengguna_id ORDER BY rental_count DESC LIMIT 5`;

  try {
    const [results] = await pool.query(query, params);
    const labels = results.map(r => r.unit_full);
    const data = results.map(r => r.rental_count);
    const agents = results.map(r => r.user_pengguna_id); // For admin view
    res.json({ success: true, labels, data, agents });
  } catch (error) {
    console.error("Popular Units Error:", error);
    res.status(500).json({ success: false, message: 'Gagal memuat unit populer.' });
  }
});


// --- 3. AGENT PERFORMANCE (ADMINS ONLY) ---
app.get('/api/dashboard/agent-performance', checkAdmin, async (req, res) => {
  const { filter } = req.query;

  let dateFilter = '';
  if (filter === 'weekly') {
    dateFilter = 'WHERE r.created_at >= DATE_SUB(NOW(), INTERVAL 7 DAY)';
  } else if (filter === 'monthly') {
    dateFilter = 'WHERE r.created_at >= DATE_SUB(NOW(), INTERVAL 1 MONTH)';
  }

  const query = `
        SELECT p.email AS user_email, COUNT(*) as total_rented 
        FROM rentals r
        LEFT JOIN Pengguna p ON r.user_pengguna_id = p.pengguna_id
        ${dateFilter} 
        GROUP BY p.email 
        ORDER BY total_rented DESC
    `;

  try {
    const [stats] = await pool.query(query);
    res.json({ success: true, data: stats });
  } catch (error) {
    console.error("Agent Performance Error:", error);
    res.status(500).json({ success: false, message: 'Gagal memuat statistik agen.' });
  }
});


// --- 4. VIOLATIONS ---
// --- 4. VIOLATIONS ---
app.post('/api/dashboard/violations', checkAdmin, upload.single('photo'), async (req, res) => {
  const { rental_id, description } = req.body;
  const photo_url = req.file ? `/uploads/${req.file.filename}` : null;
  
  try {
    // Added 'uploaded_by' to the query and passed req.user.nama
    await pool.query(
      `INSERT INTO violations (rental_id, description, photo_url, pengguna_id, uploaded_by) VALUES (?, ?, ?, ?, ?)`,
      [rental_id, description, photo_url, req.user.id, req.user.nama || req.user.email]
    );
    res.json({ success: true, message: 'Laporan pelanggaran berhasil diupload.' });
  } catch (error) {
    // Always log the error so your terminal tells you exactly what went wrong next time!
    console.error("Violation Upload Error:", error); 
    res.status(500).json({ success: false, message: 'Gagal upload laporan.' });
  }
});


app.get('/api/dashboard/violations', checkAuth, async (req, res) => {
  try {
    const user = req.user;
    const userIsAdmin = isAdmin(user);

    const uploadedByField = userIsAdmin ? 'v.pengguna_id' : "'' as pengguna_id";

    const query = `SELECT v.id, v.description, v.photo_url, v.created_at, r.nama, CONCAT(u.tower, '-', u.lantai, '-', u.unit) as unit, ${uploadedByField}
            FROM violations v 
            JOIN rentals r ON v.rental_id = r.id 
            LEFT JOIN units u ON r.unit_number = u.unit_number
            ORDER BY v.created_at DESC`;

    const [violations] = await pool.query(query);
    res.json({ success: true, data: violations });
  } catch (error) {
    console.error("Violation Fetch Error:", error);
    res.status(500).json({ success: false, message: 'Gagal memuat laporan.' });
  }
});



// Endpoint to update a violation description
app.put('/api/dashboard/violations/:id', checkAdmin, async (req, res) => {
  const { id } = req.params;
  const { description } = req.body;
  if (!description) {
    return res.status(400).json({ success: false, message: 'Deskripsi tidak boleh kosong.' });
  }
  try {
    const [result] = await pool.query('UPDATE violations SET description = ? WHERE id = ?', [description, id]);
    if (result.affectedRows === 0) {
      return res.status(404).json({ success: false, message: 'Laporan tidak ditemukan.' });
    }
    res.json({ success: true, message: 'Deskripsi laporan berhasil diperbarui.' });
  } catch (error) {
    console.error("Violation Update Error:", error);
    res.status(500).json({ success: false, message: 'Gagal memperbarui laporan.' });
  }
});

// Endpoint to delete a violation
app.delete('/api/dashboard/violations/:id', checkAdmin, async (req, res) => {
  const { id } = req.params;
  try {
    const [result] = await pool.query('DELETE FROM violations WHERE id = ?', [id]);
    if (result.affectedRows === 0) {
      return res.status(404).json({ success: false, message: 'Laporan tidak ditemukan.' });
    }
    res.json({ success: true, message: 'Laporan pelanggaran berhasil dihapus.' });
  } catch (error) {
    console.error("Violation Delete Error:", error);
    res.status(500).json({ success: false, message: 'Gagal menghapus laporan.' });
  }
});


// ================== MANAJEMEN UNIT ==================
async function getOccupiedUnits() {
  try {
    const [occupied] = await pool.query(
      `SELECT DISTINCT unit_number 
      FROM rentals 
      WHERE waktu_checkout IS NULL`
    );
    return new Set(occupied.map(r => r.unit_number));
  } catch (error) {
    console.error("Error fetching occupied units:", error);
    return new Set();
  }
}



// We map the various possible inputs (agent_email or agent_key) into the user_pengguna_id (integer)
async function resolveAgentIdFromKey(agentKey) {
  if (!agentKey) return null;

  const conn = await pool.getConnection();
  try {
    if (String(agentKey).startsWith('id:')) {
      const penggunaId = parseInt(agentKey.slice(3), 10);
      return Number.isInteger(penggunaId) ? penggunaId : null;
    }

    const [rows] = await conn.query(
      `SELECT pengguna_id FROM Pengguna WHERE email = ? LIMIT 1`,
      [agentKey]
    );
    return rows.length > 0 ? rows[0].pengguna_id : null;
  } finally {
    conn.release();
  }
}

function parseUnitNumber(unit_number) {
  const parts = (unit_number || "").split('-');
  return {
    unit_tower: parts[0] || '',
    unit_lantai: parts[1] || '',
    unit_nomor: parts[2] || '',
  };
}

app.get('/api/agents', checkAdmin, async (req, res) => {
  try {
    // Returns every Pengguna who has at least one Pelaku Komersil role (multi-role safe).
    // Uses the same join pattern confirmed working in phpMyAdmin.
    const [agents] = await pool.query(`
      SELECT DISTINCT p.pengguna_id AS id, p.nama, p.email
      FROM Pengguna p
      JOIN pengguna_role pr ON pr.pengguna_id = p.pengguna_id
      JOIN Role r ON r.role_id = pr.role_id
      WHERE LOWER(r.nama) = 'Pelaku Komersil'
      ORDER BY COALESCE(NULLIF(TRIM(p.nama), ''), p.email), p.email
    `);
    res.json({ success: true, data: agents });
  } catch (error) {
    console.error('Load Pelaku Komersil Error:', error);
    res.status(500).json({ success: false, message: 'Gagal memuat daftar pengguna Pelaku Komersil.' });
  }
});

app.get('/api/units/:agentEmail', checkAdmin, async (req, res) => {
  try {
    const agentKey = decodeURIComponent(req.params.agentEmail);
    const agentId = await resolveAgentIdFromKey(agentKey);

    if (!agentId) {
      return res.json({ success: true, data: [] });
    }

    const [units] = await pool.query("SELECT unit_number as id, unit_number FROM units WHERE user_pengguna_id = ?", [agentId]);
    const formattedUnits = units.map(u => ({ id: u.id, unit_number: u.unit_number, ...parseUnitNumber(u.unit_number) }));
    res.json({ success: true, data: formattedUnits });
  } catch (error) {
    res.status(500).json({ success: false, message: 'Gagal memuat unit agen.' });
  }
});

app.get('/api/units', checkAuth, async (req, res) => {
  try {
    const [units] = await pool.query(
      `SELECT unit_number as unit_id,unit_number, tower, lantai, unit as nomor_unit, null as kode 
             FROM units ORDER BY tower, lantai, unit`
    );

    const occupiedSet = await getOccupiedUnits();

    const data = units.map(u => {
      const unitString = `${u.tower}-${u.lantai}-${u.nomor_unit}`;
      return {
        ...u,
        status: occupiedSet.has(unitString) ? 'Occupied' : 'Available'
      };
    });

    res.json({ success: true, data });
  } catch (error) {
    console.error("Error loading units:", error);
    res.status(500).json({ success: false, message: "Gagal memuat unit." });
  }
});


app.post('/api/units', checkAdmin, async (req, res) => {
  const { unit_tower, unit_lantai, unit_nomor, agent_email, agent_key } = req.body;

  let resolvedAgentId = await resolveAgentIdFromKey(agent_email);
  if (!resolvedAgentId && agent_key) {
    resolvedAgentId = await resolveAgentIdFromKey(agent_key);
  }

  if (!unit_tower || !unit_lantai || !unit_nomor || !resolvedAgentId) {
    return res.status(400).json({ success: false, message: "Data unit tidak lengkap." });
  }

  try {
    const unit_number = `${unit_tower}-${unit_lantai}-${unit_nomor}`;
    await pool.query(
      `INSERT INTO units (unit_number, user_pengguna_id, tower, lantai, unit) VALUES (?, ?, ?, ?, ?)`,
      [unit_number, resolvedAgentId, unit_tower, unit_lantai, unit_nomor]
    );
    res.json({ success: true, message: "Unit berhasil ditambahkan." });
  } catch (error) {
    console.error("Add Unit Error:", error);
    res.status(500).json({ success: false, message: "Gagal menambah unit." });
  }
});

app.put('/api/units/:id', checkAdmin, async (req, res) => {
  const { id } = req.params;
  const { tower, lantai, nomor_unit, kode } = req.body;

  try {
    const [result] = await pool.query(
      `UPDATE units SET tower = ?, lantai = ?, unit = ?, unit_number = ? WHERE unit_number = ?`,
      [tower, lantai, nomor_unit, `${tower}-${lantai}-${nomor_unit}`, id]
    );

    if (result.affectedRows === 0) {
      return res.status(404).json({ success: false, message: "Unit tidak ditemukan." });
    }

    res.json({ success: true, message: "Unit berhasil diperbarui." });
  } catch (error) {
    console.error("Update Unit Error:", error);
    res.status(500).json({ success: false, message: "Gagal memperbarui unit." });
  }
});


app.delete('/api/units/:id', checkAdmin, async (req, res) => {
  const { id } = req.params;

  try {
    const [inUse] = await pool.query(
      `SELECT id FROM rentals 
             WHERE unit_number = (
                SELECT unit_number 
                FROM units WHERE unit_number = ?
             ) AND waktu_checkout IS NULL`,
      [id]
    );

    if (inUse.length > 0) {
      return res.status(400).json({
        success: false,
        message: "Unit tidak dapat dihapus karena sedang digunakan."
      });
    }

    const [result] = await pool.query(`DELETE FROM units WHERE unit_number = ?`, [id]);

    if (result.affectedRows === 0) {
      return res.status(404).json({ success: false, message: "Unit tidak ditemukan." });
    }

    res.json({ success: true, message: "Unit berhasil dihapus." });
  } catch (error) {
    console.error("Delete Unit Error:", error);
    res.status(500).json({ success: false, message: "Gagal menghapus unit." });
  }
});


app.get('/api/my-units', checkAuth, async (req, res) => {
  try {
    const occupiedUnitsSet = await getOccupiedUnits();

    const [units] = await pool.query("SELECT unit_number as id, unit_number FROM units WHERE user_pengguna_id = ?", [req.user.id]);
    const formattedUnits = units.map(u => {

      const unitParts = parseUnitNumber(u.unit_number);
      const unitString = `${unitParts.unit_tower}-${unitParts.unit_lantai}-${unitParts.unit_nomor}`;
      return {
        id: u.id,
        unit_number: u.unit_number,
        ...unitParts,
        is_occupied: occupiedUnitsSet.has(unitString)
      };

    });
    res.json({ success: true, data: formattedUnits });
  } catch (error) {
    res.status(500).json({ success: false, message: 'Gagal memuat unit Anda.' });
  }
});


// Get all units for admin roles
app.get('/api/all-units', checkAdmin, async (req, res) => {
  try {
    const occupiedUnitsSet = await getOccupiedUnits();

    const query = 'SELECT u.unit_number as id, u.unit_number, u.user_pengguna_id as agent_email ' +
      'FROM units u ' +
      'ORDER BY u.user_pengguna_id, u.unit_number';

    const [units] = await pool.query(query);
    const formattedUnits = units.map(u => {
      const unitParts = parseUnitNumber(u.unit_number);
      const unitString = `${unitParts.unit_tower}-${unitParts.unit_lantai}-${unitParts.unit_nomor}`;
      return {
        id: u.id,
        unit_number: u.unit_number,
        ...unitParts,
        agent_email: u.agent_email,
        is_occupied: occupiedUnitsSet.has(unitString)
      };
    });
    res.json({ success: true, data: formattedUnits });
  } catch (error) {
    console.error("Error fetching all units:", error);
    res.status(500).json({ success: false, message: 'Gagal memuat semua unit.' });
  }
});

// Helper function to calculate duration
function calculateDays(checkin, checkout) {
  const d1 = new Date(checkin);
  const d2 = new Date(checkout);
  if (isNaN(d1) || isNaN(d2)) return 0;
  const diffTime = Math.max(0, d2 - d1);
  return Math.ceil(diffTime / (1000 * 60 * 60 * 24));
}

function parseHourValue(value, fallback = 0) {
  if (!value) return fallback;
  const [hours] = String(value).split(':');
  const parsedHour = parseInt(hours, 10);
  return Number.isFinite(parsedHour) ? parsedHour : fallback;
}

function predictRentalClassification({ lama_menginap, waktu_checkin, waktu_checkout, jenis_sewa, status_pasutri, status_kewarganegaraan, metode_pembayaran }) {
  const payload = {
    lama_menginap: Number.parseInt(lama_menginap, 10) || 0,
    checkin_hour: parseHourValue(waktu_checkin, 0),
    checkout_hour: parseHourValue(waktu_checkout || waktu_checkin, parseHourValue(waktu_checkin, 0)),
    jenis_sewa: jenis_sewa || 'Unknown',
    status_pasutri: status_pasutri || 'Unknown',
    status_kewarganegaraan: status_kewarganegaraan || 'Unknown',
    metode_pembayaran: metode_pembayaran || 'Unknown'
  };

  const pythonCode = `
import json
import os
import warnings
from pathlib import Path

import joblib
import pandas as pd

warnings.filterwarnings('ignore')

payload = json.loads(os.environ['PREDICT_PAYLOAD'])
bundle = joblib.load(Path(os.environ['MODEL_PATH']))
model = bundle['model']
label_encoders = bundle['label_encoders']
selected_features = bundle['selected_features']

def encode_value(feature_name, value):
    encoder = label_encoders.get(feature_name)
    if encoder is None:
        return value
    candidate = value if value in list(encoder.classes_) else 'Unknown'
    return int(encoder.transform([candidate])[0])

row = {
    'lama_menginap': int(payload.get('lama_menginap') or 0),
    'checkin_hour': int(payload.get('checkin_hour') or 0),
    'checkout_hour': int(payload.get('checkout_hour') or 0),
    'jenis_sewa': encode_value('jenis_sewa', payload.get('jenis_sewa', 'Unknown')),
    'status_pasutri': encode_value('status_pasutri', payload.get('status_pasutri', 'Unknown')),
    'status_kewarganegaraan': encode_value('status_kewarganegaraan', payload.get('status_kewarganegaraan', 'Unknown')),
    'metode_pembayaran': encode_value('metode_pembayaran', payload.get('metode_pembayaran', 'Unknown'))
}

# Safely build the feature array, defaulting to 0 if a feature is somehow missing to prevent NaN errors
final_features = [row.get(feature, 0) for feature in selected_features]
frame = pd.DataFrame([final_features], columns=selected_features)

prediction = int(model.predict(frame)[0])
probability = None

if hasattr(model, 'predict_proba'):
    probability = float(model.predict_proba(frame)[0][prediction])

print(json.dumps({
    'class': prediction,
    'label': 'Aman' if prediction == 0 else 'Melanggar',
    'notification': '??Penyewa Berpotensi Melakukan Pelanggaran!' if prediction == 1 else None,
    'probability': probability
}))
`;

  const result = spawnSync(PYTHON_EXECUTABLE, ['-c', pythonCode], {
    env: {
      ...process.env,
      MODEL_PATH: MODEL_PATH,
      PREDICT_PAYLOAD: JSON.stringify(payload)
    },
    encoding: 'utf8',
    maxBuffer: 10 * 1024 * 1024
  });

  if (result.error) {
    throw result.error;
  }

  if (result.status !== 0) {
    throw new Error(result.stderr || 'Gagal menjalankan prediksi model.');
  }

  return JSON.parse((result.stdout || '').trim() || '{}');
}

// NEW ENDPOINT: Save a new rental record
app.post('/api/rentals', checkAuth, async (req, res) => {
  const conn = await pool.getConnection();
  try {
    const {
      nama, nik, status_pasutri, unit_number, status_kewarganegaraan,
      jenis_sewa, metode_pembayaran, metode_lain, tanggal_checkin, waktu_checkin, tanggal_checkout
    } = req.body;

    const normalizedNik = String(nik || '').replace(/\D/g, '');
    if (!/^\d{16}$/.test(normalizedNik)) {
      return res.status(400).json({ success: false, message: 'NIK harus terdiri dari 16 digit angka.' });
    }

    const userIsAdmin = isAdmin(req.user);
    let agent_id = req.user.id;
    if (userIsAdmin && req.body.agent_id) {
      agent_id = req.body.agent_id;
    }

    const lama_menginap = calculateDays(tanggal_checkin, tanggal_checkout);
    const db_status_pasutri = status_pasutri === 'Ya' ? 'Menikah' : 'Belum Menikah';

    const sql = `INSERT INTO rentals (
            nama, nik, status_pasutri, status_kewarganegaraan, 
            jenis_sewa, unit_number, metode_pembayaran, metode_lain, tanggal_checkin, waktu_checkin, 
            tanggal_checkout, lama_menginap, user_pengguna_id
        ) VALUES (?, ?, ?, ?, ?, ?, ?, ?, ?, ?, ?, ?, ?)`;

    const params = [
      nama, normalizedNik, db_status_pasutri, status_kewarganegaraan,
      jenis_sewa, unit_number, metode_pembayaran, metode_pembayaran === 'Others' ? metode_lain : null,
      tanggal_checkin, waktu_checkin, tanggal_checkout, lama_menginap, agent_id
    ];

    await conn.beginTransaction();
    const [result] = await conn.query(sql, params);
    const newRentalId = result.insertId;

    // Add creation event to logs
    const logQuery = `INSERT INTO rental_logs (action, field_changed, new_value, pengguna_id, rental_id) VALUES ?`;
    const logEntries = [
      ['create', 'all', JSON.stringify(req.body), req.user.id, newRentalId]
    ];
    await conn.query(logQuery, [logEntries]);

    await conn.commit();
    
    let prediction = null;
    try {
      prediction = predictRentalClassification({
        lama_menginap,
        waktu_checkin,
        waktu_checkout: req.body.waktu_checkout,
        jenis_sewa,
        status_pasutri: db_status_pasutri,
        status_kewarganegaraan,
        metode_pembayaran
      });
    } catch (predictionError) {
      console.error('Prediction Error:', predictionError);
    }

    res.json({
      success: true,
      message: "Data penyewa berhasil disimpan.",
      prediction
    });

  } catch (error) {
    if (conn) await conn.rollback();
    console.error("Error saving rental:", error);
    res.status(500).json({ success: false, message: 'Gagal menyimpan data penyewa.' });
  } finally {
    if (conn) conn.release();
  }
});

// ================== ETC ==================
app.get('/api/pelanggaran-list', checkAuth, (req, res) => {
  const user = req.user;
  const listForAgen = [
    "Ditemukan alat suntik di tempat sampah", "Ditemukan kondom dalam jumlah banyak",
    "Kerusakan parah pada fasilitas", "Kebisingan berlebihan di malam hari",
    "Penyalahgunaan alkohol/narkoba", "Kekerasan atau ancaman kepada penghuni lain",
    "Merokok di area terlarang", "Tidak menjaga kebersihan unit",
    "Terpantau adanya tamu yang keluar masuk pada malam hari", "Menyewakan kembali unit yang disewa",
  ];
  const listForAdmin = [
    ...listForAgen,
    "Agen memberlakukan sistem transit", "Agen lalai terhadap pelanggaran penyewa"
  ];

  if (isAdmin(user)) {
    res.json({ success: true, data: listForAdmin });
  } else {
    res.json({ success: true, data: listForAgen });
  }
});

// Login page (served separately to avoid redirect loop)
app.get('/login', (req, res) => {
  res.sendFile(path.join(__dirname, 'public', 'login.html'));
});

// Fallback Route
app.get('*', (req, res) => {
  res.sendFile(path.join(__dirname, 'public', 'index.html'));
});

// Server Start
app.listen(port, () => {
  console.log(`? Server berjalan di http://localhost:${port}`);
});
\end{lstlisting}

\begin{lstlisting}[language=JavaScript, caption={fix.js}]
const fs = require('fs');
let code = fs.readFileSync('app.js', 'utf8');
const start = code.indexOf("app.post('/api/login',");
const end = code.indexOf("app.post('/api/logout',");

if (start > -1 && end > -1) {
  const newLogin = app.post('/api/login', async (req, res) => {
  const { email, password } = req.body;
  try {
    const [rows] = await pool.query('SELECT * FROM users WHERE email = ?', [email]);
    if (rows.length === 0) {
      return res.status(401).json({ success: false, message: 'Email tidak ditemukan.' });
    }
    const user = rows[0];
    const bcrypt = require('bcryptjs');
    const isMatch = await bcrypt.compare(password, user.password);
    if (!isMatch) {
      return res.status(401).json({ success: false, message: 'Password salah.' });
    }
    const roles = [user.role];
    const payload = { id: user.id, email: user.email, nama: user.name, nib: user.nib, roles };
    const token = jwt.sign(payload, JWT_SECRET, { expiresIn: '12h' });
    res.json({ success: true, token, user: payload });
  } catch (error) {
    console.error('Login Error:', error);
    res.status(500).json({ success: false, message: 'Server error.' });
  }
});

app.post('/api/register', async (req, res) => {
  const { name, email, password, nib } = req.body;
  if (!email || !password || !nib || !name) {
    return res.status(400).json({ success: false, message: 'Semua kolom harus diisi.' });
  }
  if (nib.length !== 13) {
    return res.status(400).json({ success: false, message: 'NIB harus 13 digit.' });
  }
  try {
    const [existing] = await pool.query('SELECT * FROM users WHERE email = ?', [email]);
    if (existing.length > 0) {
      return res.status(400).json({ success: false, message: 'Email sudah terdaftar.' });
    }
    const bcrypt = require('bcryptjs');
    const hashedPassword = await bcrypt.hash(password, 10);
    await pool.query('INSERT INTO users (email, password, nib, role, name) VALUES (?, ?, ?, ?, ?)', [email, hashedPassword, nib, 'agen', name]);
    res.json({ success: true, message: 'Registrasi berhasil, silakan login.' });
  } catch (error) {
    console.error('Register Error:', error);
    res.status(500).json({ success: false, message: 'Server error.' });
  }
});

;
  code = code.slice(0, start) + newLogin + code.slice(end);
  fs.writeFileSync('app.js', code);
  console.log('Replaced');
} else {
  console.log('Not found');
}
\end{lstlisting}


\begin{lstlisting}[language=JavaScript, caption={findindex.js}]
const fs = require('fs');
let code = fs.readFileSync('public/index.html', 'utf8');

code = code.replace(/window\.location\.href = 'https\:\/\/member\.thejarrdin\.com\/login\?from=visitor';/g, "window.location.href = '/login.html'; // window.location.href = 'https://member.thejarrdin.com/login?from=visitor';");
code = code.replace(/window\.location\.href = 'https\:\/\/member\.thejarrdin\.com\/logout';/g, "window.location.href = '/login.html'; // window.location.href = 'https://member.thejarrdin.com/logout';");

fs.writeFileSync('public/index.html', code);
console.log('Replaced in index.html');
\end{lstlisting}


\begin{lstlisting}[language=JavaScript, caption={fixlogin.js}]
const fs = require('fs');
let code = fs.readFileSync('app.js', 'utf8');
const start = code.indexOf("app.post('/api/login',");
const end = code.indexOf("app.post('/api/logout',");

if (start > -1 && end > -1) {
  const newLogin = `app.post('/api/login', async (req, res) => {
  const { email, password } = req.body;
  try {
    const [rows] = await pool.query('SELECT * FROM users WHERE email = ?', [email]);
    if (rows.length === 0) {
      return res.status(401).json({ success: false, message: 'Email tidak ditemukan.' });
    }
    const user = rows[0];
    const bcrypt = require('bcryptjs');
    const isMatch = await bcrypt.compare(password, user.password);
    if (!isMatch) {
      return res.status(401).json({ success: false, message: 'Password salah.' });
    }
    const roles = [user.role];
    const payload = { id: user.id, email: user.email, nama: user.name, nib: user.nib, roles };
    const token = jwt.sign(payload, JWT_SECRET, { expiresIn: '12h' });
    res.json({ success: true, token, user: payload });
  } catch (error) {
    console.error('Login Error:', error);
    res.status(500).json({ success: false, message: 'Server error.' });
  }
});

app.post('/api/register', async (req, res) => {
  const { name, email, password, nib } = req.body;
  if (!email || !password || !nib || !name) {
    return res.status(400).json({ success: false, message: 'Semua kolom harus diisi.' });
  }
  if (nib.length !== 13) {
    return res.status(400).json({ success: false, message: 'NIB harus 13 digit.' });
  }
  try {
    const [existing] = await pool.query('SELECT * FROM users WHERE email = ?', [email]);
    if (existing.length > 0) {
      return res.status(400).json({ success: false, message: 'Email sudah terdaftar.' });
    }
    const bcrypt = require('bcryptjs');
    const hashedPassword = await bcrypt.hash(password, 10);
    await pool.query('INSERT INTO users (email, password, nib, role, name) VALUES (?, ?, ?, ?, ?)', [email, hashedPassword, nib, 'agen', name]);
    res.json({ success: true, message: 'Registrasi berhasil, silakan login.' });
  } catch (error) {
    console.error('Register Error:', error);
    res.status(500).json({ success: false, message: 'Server error.' });
  }
});

`;
  code = code.slice(0, start) + newLogin + code.slice(end);
  fs.writeFileSync('app.js', code);
  console.log('Replaced');
} else {
  console.log('Not found');
}
\end{lstlisting}



























\begin{lstlisting}[language=Python, caption={model-api.py}]
import os
import re
import warnings
from pathlib import Path

from flask import Flask, request, jsonify
from flask_cors import CORS
import jwt
import joblib
import pandas as pd
from dotenv import load_dotenv

warnings.filterwarnings('ignore')

# Load .env from current directory for standalone deployment, fallback to parent for local mono-repo
env_path = Path(__file__).resolve().parent / '.env'
if not env_path.exists():
    env_path = Path(__file__).resolve().parent.parent / '.env'
load_dotenv(dotenv_path=env_path)

app = Flask(__name__)
CORS(app)

JWT_SECRET = os.environ.get('JWT_SECRET', 'jwt-secret-key-jarrdin')
SRUSUN_JWT_SECRET = os.environ.get('SRUSUN_JWT_SECRET', 'jarrdin-cihampelas')

# Pre-load the model to save time on each request
MODEL_PATH = Path(__file__).resolve().parent.parent / 'model_xgboost.pkl'
try:
    bundle = joblib.load(MODEL_PATH)
    model = bundle['model']
    label_encoders = bundle['label_encoders']
    selected_features = bundle['selected_features']
    print(f"Model loaded successfully from {MODEL_PATH}")
except Exception as e:
    print(f"Error loading model from {MODEL_PATH}: {e}")
    model, label_encoders, selected_features = None, None, []

def check_auth(token):
    if not token:
        return False, "Akses ditolak. Silakan login."
    
    try:
        jwt.decode(token, JWT_SECRET, algorithms=["HS256"])
        return True, None
    except Exception:
        try:
            jwt.decode(token, SRUSUN_JWT_SECRET, algorithms=["HS256"])
            return True, None
        except Exception:
            return False, "Token tidak valid. Silakan login ulang."

def encode_value(feature_name, value):
    if label_encoders is None:
        return value
    encoder = label_encoders.get(feature_name)
    if encoder is None:
        return value
    candidate = value if value in list(encoder.classes_) else 'Unknown'
    return int(encoder.transform([candidate])[0])

def parse_hour_value(time_str, default_hour=0):
    if not time_str:
        return default_hour
    match = re.match(r'^(\d{1,2}):', str(time_str))
    if match:
        return int(match.group(1))
    return default_hour

@app.route('/api/predict', methods=['POST'])
def predict():
    auth_header = request.headers.get('Authorization')
    token = None
    if auth_header and auth_header.startswith('Bearer '):
        token = auth_header.split(' ')[1]
    else:
        token = request.args.get('token')

    is_valid, error_msg = check_auth(token)
    if not is_valid:
        return jsonify({'success': False, 'message': error_msg}), 401

    if model is None:
        return jsonify({'success': False, 'message': 'Model not loaded properly.'}), 500

    data = request.json or {}
    
    try:
        lama_menginap = int(data.get('lama_menginap') or 0)
        
        waktu_checkin = data.get('waktu_checkin')
        waktu_checkout = data.get('waktu_checkout')
        
        checkin_hour = parse_hour_value(waktu_checkin, 0)
        checkout_hour = parse_hour_value(waktu_checkout or waktu_checkin, parse_hour_value(waktu_checkin, 0))
        
        row = {
            'lama_menginap': lama_menginap,
            'checkin_hour': checkin_hour,
            'checkout_hour': checkout_hour,
            'jenis_sewa': encode_value('jenis_sewa', data.get('jenis_sewa', 'Unknown')),
            'status_pasutri': encode_value('status_pasutri', data.get('status_pasutri', 'Unknown')),
            'status_kewarganegaraan': encode_value('status_kewarganegaraan', data.get('status_kewarganegaraan', 'Unknown')),
            'metode_pembayaran': encode_value('metode_pembayaran', data.get('metode_pembayaran', 'Unknown'))
        }

        final_features = [row.get(feature, 0) for feature in selected_features]
        frame = pd.DataFrame([final_features], columns=selected_features)

        prediction = int(model.predict(frame)[0])
        probability = None

        if hasattr(model, 'predict_proba'):
            probability = float(model.predict_proba(frame)[0][prediction])

        return jsonify({
            'success': True,
            'prediction': {
                'class': prediction,
                'label': 'Aman' if prediction == 0 else 'Melanggar',
                'notification': 'Penyewa Berpotensi Melakukan Pelanggaran!' if prediction == 1 else None,
                'probability': probability
            }
        })
    except Exception as e:
        print("Prediction API Error:", e)
        return jsonify({'success': False, 'message': 'Prediction API Error', 'error': str(e)}), 500

if __name__ == '__main__':
    port = int(os.environ.get('MODEL_API_PORT', 3001))
    app.run(host='0.0.0.0', port=port, debug=False)

\end{lstlisting}

